%% LyX 2.4.0-beta3-devel created this file.  For more info, see https://www.lyx.org/.
%% Do not edit unless you really know what you are doing.
\documentclass[english,tableposition=top]{report}
\usepackage{lmodern}
\renewcommand{\sfdefault}{lmss}
\renewcommand{\ttdefault}{lmtt}
\usepackage[T1]{fontenc}
\usepackage{textcomp}
\usepackage[latin9]{inputenc}
\setcounter{secnumdepth}{3}
\usepackage{color}
\definecolor{shadecolor}{rgb}{0.667969, 1, 1}
\usepackage{babel}
\usepackage{wrapfig}
\usepackage{booktabs}
\usepackage{framed}
\usepackage{url}
\usepackage{amsmath}
\usepackage{graphicx}
\usepackage[pdfusetitle,
 bookmarks=true,bookmarksnumbered=true,bookmarksopen=true,bookmarksopenlevel=2,
 breaklinks=true,pdfborder={0 0 1},backref=section,colorlinks=true]
 {hyperref}

\makeatletter

%%%%%%%%%%%%%%%%%%%%%%%%%%%%%% LyX specific LaTeX commands.
%% Because html converters don't know tabularnewline
\providecommand{\tabularnewline}{\\}

%%%%%%%%%%%%%%%%%%%%%%%%%%%%%% Textclass specific LaTeX commands.
\newenvironment{centred}%
  {\begin{center}}{\end{center}}

%%%%%%%%%%%%%%%%%%%%%%%%%%%%%% User specified LaTeX commands.
\usepackage{numerica}
\usepackage{numerica-plus}
\usepackage{upquote}
\newcommand\rel{\,\varrho\;}
\DeclareMathOperator{\erf}{erf}
\DeclareMathOperator{\gd}{gd}

\reuse{}
\constants{ c=30, \omega=0.2 }

\makeatother

\begin{document}
\title{\texttt{numerica-plus}~\\
}
\author{Andrew Parsloe\\
(\url{ajparsloe@gmail.com})}
\maketitle
\begin{abstract}
The \verb`numerica-plus` package defines commands to iterate functions
of a single variable, find fixed points of such functions, find the
zeros or extrema of such functions, and calculate the terms of recurrence
relations.
\end{abstract}
\noindent\begin{minipage}[t]{1\columnwidth}%
\begin{shaded}%

\paragraph*{Note:}
\begin{itemize}
\item This document applies to version 3.0.0 of \texttt{numerica-plus.}
\item Version 3 of \texttt{numerica} is required; (\texttt{numerica} requires
\texttt{amsmath} and \texttt{mathtools}).
\item I refer many times in this document to \emph{Handbook of Mathematical
Functions}, edited by Milton Abramowitz and Irene A. Stegun, Dover,
1965. This is abbreviated to \emph{HMF}, and often followed by a number
like 1.2.3 to locate the actual expression or value referenced.
\item Version 3.0.0 of \texttt{numerica-plus}
\begin{itemize}
\item is compatible with the additional features of numerica version 3.0.0,
\item including the decimal comma if the \verb`comma` package option is
used with numerica;
\item amends and adds to documentation, including
\item a section on finding roots with Newton-Raphson iteration.
\end{itemize}
\end{itemize}
\end{shaded}%
\end{minipage}

\tableofcontents{}

\chapter{Introduction}

Entering
\begin{verbatim}
    \usepackage[<options>]{numerica}
    \usepackage{numerica-plus}
\end{verbatim}
in the preamble of your document makes available the commands
\begin{itemize}
\item \verb`\nmcIterate`, a command to iterate a function (apply it repeatedly
to itself), including finding fixed points (values of $x$ where $f(x)=x$);
\item \verb`\nmcSolve`, a command to find the zeros of functions of a single
variable (values of $x$ for which $f(x)=0$) or, failing that, local
maxima or minima of such functions;
\item \verb`\nmcRecur`, a command to calculate the values of terms in recurrence
relations in a single recurrence variable (like the terms of the Fibonacci
sequence or Legendre polynomials).
\end{itemize}
A main difference from version 2 of \verb`numerica-plus` is that
the package now does not load \texttt{numerica} automatically but,
rather, requires the user to explicitly call \texttt{numerica} (with
options if wanted). Version 3 of \texttt{numerica} is required, and
must be loaded \emph{before} \texttt{numerica-plus}. With \texttt{numerica}
loaded you get access to the commands \verb`\nmcEvaluate` (\verb`\eval`),
\verb`\nmcInfo` (\verb`\info`), \verb`\nmcMacros` (\verb`\macros`),
\verb`\nmcConstants` (\verb`\constants`), and \verb`\nmcReuse`
(\verb`\reuse`); see the \texttt{numerica} documentation for details
on the use of these commands. The \texttt{numerica} package options,
and particularly the \verb`comma` and \verb`rounding` options specifying
the decimal comma and default rounding value, apply in \texttt{numerica-plus}.

The commands of the present package all share the syntax of \verb`\nmcEvaluate`.
I will discuss them individually in later chapters but turn first
to a meaningful example to illustrate their use and give a sense of
`what they are about'.

\section{Example of use: the rotating disk}

\label{sec:introExampleOfUse}Consider a disk rotating uniformly with
angular velocity $\omega$ in an anticlockwise sense in an inertial
system in which the disk's centre \textbf{0} is at rest. Three distinct
points \textbf{1}, \textbf{2}, \textbf{3} are fixed in the disk and,
in a co-rotating polar coordinate system centred at \textbf{0}, have
polar coordinates $(r_{i},\theta_{i})$ ($i,j=1,2,3$). Choose \textbf{01}
as initial line so that $\theta_{1}=0$.

The cosine rule for solving triangles tells us that the time $t_{ij}$
in the underlying inertial system for a signal to pass from point
\textbf{i} to point \textbf{j} satisfies the equation
\[
t_{ij}=c^{-1}\sqrt{r_{i}^{2}+r_{j}^{2}-2r_{i}r_{j}\cos(\theta_{j}-\theta_{i}+\omega t_{ij})}\equiv f(t_{ij}),
\]
where $c$ is the speed of light and $i,j\in\{1,2,3\}$. (Equally,
we could be describing an acoustic signal between points on a disk
rotating uniformly in a still, uniform atmosphere \textendash{} in
which case $c$ would be the speed of sound.) Although the equation
doesn't solve algebraically for the time $t_{ij},$ it does tell us
that $t=t_{ij}$ is a \emph{fixed point} of the function $f$. To
calculate fixed points we use the command \verb`\nmcIterate`, or
its short-name form \verb`\iter`, with the star option, \verb`\iter*`.
For \verb`\iter` the star option means: continue iterating until
a fixed point has been reached and, as with the \verb`\eval` command,
suppress all elements from the display save for the numerical result. 

First, though, values need to be assigned to the various parameters.
Suppose we use units in which $c=30,$ and $\omega=0.2$ radians per
second. To avoid having to write these values in the vv-list every
time, I have put in the preamble to this document the statement
\begin{verbatim}
    \constants{ c=30,\omega=0.2 }
\end{verbatim}
For the polar coordinates of \textbf{1 }and \textbf{3 }I have chosen
$r_{1}=10$, $r_{3}=20$ and $\theta_{3}=0.2$ radians (remember $\theta_{1}=0$).
To find a fixed point $t_{13}$ I give $t$ an initial trial value
$1$ (plucked from the air). Its position as the \emph{rightmost}
item in the vv-list tells \verb`\iter` that $t$ is the iteration
variable:
\begin{verbatim}
    \iter*{ c^{-1}\sqrt{r_1^2+r_3^2-2r_1 r_3
              \cos(\theta_3+\omega t)}
          }[ r_1=10,r_3=20,\theta_3=0.2,t=1 ],
    \quad\info{iter}.
\end{verbatim}
$\Longrightarrow$     \iter*{ c^{-1}\sqrt{r_1^2+r_3^2-2r_1 r_3
              \cos(\theta_3+\omega t)}
          }[ r_1=10,r_3=20,\theta_3=0.2,t=1 ],
    \quad\info{iter}. The short-name form of the \verb`\nmcInfo` command from \verb`numerica`
 has been used to display the number of iterations required to attain
the fixed-point value. 

To six figures, only five iterations are needed, which seems rapid
but we can check this by substituting $t=0.356899$ back into the
formula and \verb`\eval`-uating it:
\begin{verbatim}
    \eval*{ c^{-1}\sqrt{r_1^2+r_3^2-2r_1 r_3
              \cos(\theta_3+\omega t)}
          }[ r_1=10,r_3=20,\theta_3=0.2,t=0.356899 ]
\end{verbatim}
$\Longrightarrow$     \eval*{ c^{-1}\sqrt{r_1^2+r_3^2-2r_1 r_3
              \cos(\theta_3+\omega t)}
          }[ r_1=10,r_3=20,\theta_3=0.2,t=0.356899 ], confirming that we have indeed calculated a fixed point. That it
did indeed take only $5$ iterations can be checked by omitting the
asterisk from the \verb`\iter` command and specifying the total number
of iterations to perform. I choose \texttt{do=}7 to show not just
the $5$th iteration but also the next two just to confirm that the
result is stable. We shall view all $7$: \texttt{see=7}. Because
of the length of the formula I have suppressed display of the vv-list
by giving the key \texttt{vv}\emph{ }an empty value:
\begin{verbatim}
    \iter[do=7,see=7,vv=]
      {\[ c^{-1}\sqrt{r_1^2+r_3^2-2r_1 r_3
        \cos(\theta_3+\omega t)} \]}
          [ r_1=10,r_3=20,\theta_3=0.2,t=1 ]
\end{verbatim}
$\Longrightarrow$     \iter[do=7,see=7,vv=]
      {\[ c^{-1}\sqrt{r_1^2+r_3^2-2r_1 r_3
        \cos(\theta_3+\omega t)}\]}
          [ r_1=10,r_3=20,\theta_3=0.2,t=1 ]
\noindent\begin{flushleft}
The display makes clear that on the $5$th iteration, the $6$-figure
value has been attained. 
\par\end{flushleft}

Alternatively, we could use the \verb`\nmcRecur` command (or its
short-name form \verb`\recur`) to view the successive iterations,
since an iteration is a first-order recurrence: $f_{n+1}=f(f_{n})$:
\begin{verbatim}
    \recur[env=multline*,vv={,}\\(vv)\\,
                   do=8,see1=0,see2=5]
      { f_{n+1}=c^{-1}\sqrt{r_1^2+r_3^2-2r_1 r_3
        \cos(\theta_3+\omega f_{n})} }
          [ r_1=10,r_3=20,\theta_3=0.2,f_{0}=1 ]
\end{verbatim}
$\Longrightarrow$     \recur[env=multline*,vv={,}\\(vv)\\,
                   do=8,see1=0,see2=5]
      { f_{n+1}=c^{-1}\sqrt{r_1^2+r_3^2-2r_1 r_3
        \cos(\theta_3+\omega f_{n})} }
          [ r_1=10,r_3=20,\theta_3=0.2,f_{0}=1 ]

\noindent I have specified \texttt{do=8} terms rather than $7$ since
the zero-th term ($f_{0}=1$) is included in the count. I've chosen
to view the last $5$ of them but none prior to those by writing \texttt{see1=0,see2=5}.
Note the choice of environment and the \texttt{vv} setting, pushing
display of the vv-list and result to new lines and suppressing equation
numbering with the \texttt{{*}} on the \verb`multline`.

Another and perhaps more obvious way to find the value of $t_{13}$,
is to look for a zero of the function $f(t)-t$. That means using
the command \verb`\nmcSolve` (or its short-name form \verb`\solve`).
I shall do so with the star option \verb`\solve*` which suppresses
display of all but the numerical result. A trial value for $t$ is
required. I have chosen \texttt{t=0}:
\begin{verbatim}
    \solve*{ c^{-1}\sqrt{r_1^2+r_3^2-2r_1 r_3
      \cos(\theta_3+\omega t)} - t }
        [ r_1=10,r_3=20,\theta_3=0.2,t=0 ],
    \quad \nmcInfo{solve}.
\end{verbatim}
$\Longrightarrow$     \solve*{ c^{-1}\sqrt{r_1^2+r_3^2-2r_1 r_3
      \cos(\theta_3+\omega t)} - t }
        [ r_1=10,r_3=20,\theta_3=0.2,t=0 ],
    \quad\nmcInfo{solve}. 

Nearly the same answer as before is attained but this time many more
steps have been required. This is to be expected. The \verb`\solve`
command uses the bisection method, finding an interval where the function
has opposite signs at the end points and successively bisecting it
to locate the zero between. Since $1/2^{10}\approx1/10^{3}$, about
$10$ bisections are needed to determine $3$ decimal places. Hence
we can expect about $20$ bisections for a $6$-decimal-place answer.
The particular form of the \verb`\nmcInfo` command display, `$1+20$
steps', indicates that it took $1$ search step to find an interval
where the function had opposite signs at the end points and, within
that interval, $20$ bisections to narrow the position of the zero
to $6$-figures. I will discuss the discrepancy in the final figure
in Chapter~\ref{chap:solveSolve}; see �\ref{subsec:solveExtraRounding}.

\subsection{Circuits}

Okay, so we can calculate the time taken in the underlying inertial
system for a signal to pass from one point of the rotating disk to
another. How long does it take to traverse the circuit \textbf{1231},
i.e. a signal from \textbf{1} to \textbf{2} to \textbf{3} and back
to \textbf{1}? That means forming the sum $t_{12}+t_{23}+t_{31}$,
which means calculating the separate $t_{ij}$ and then using \verb`\eval`
to calculate their sum. 

To simplify things, I assume a little symmetry. Let the (polar) coordinates
of \textbf{1} be $(a,0),$ of \textbf{2} be $(r,-\theta)$, and of
\textbf{3} be $(r,\theta)$: \textbf{2} and \textbf{3} are at the
same radial distance from the centre \textbf{0} and at the same angular
distance from the line \textbf{01} but on opposite sides of it, \textbf{3}
ahead of the line, \textbf{2} behind it. The rotation is in the direction
of positive $\theta$. Rather than just calculate $t_{12}+t_{23}+t_{31}$
for the circuit \textbf{1231}, I also calculate the time $t_{13}+t_{32}+t_{21}$
for a signal to traverse the same circuit but in the opposite sense,
\textbf{1321},  and compare them (form the difference). 

Note that with \textbf{2} and \textbf{3} positioned as they are relative
to \textbf{1}, a signal against the rotation from \textbf{3} to \textbf{1}
takes the same time as a signal from \textbf{1} to \textbf{2} and,
in the sense of rotation, a signal from \textbf{2} to \textbf{1} takes
the same time as a signal from \textbf{1} to \textbf{3}, so that the
round trip times are $2t_{12}+t_{23}$ and $2t_{13}+t_{32}$. 

\noindent{}%
\noindent\begin{minipage}[t]{1\columnwidth}%
\begin{shaded}%
To see this, suppose the signal from \textbf{2} to \textbf{1} starts
at time $t=0$; it reaches \textbf{1} at a later time $t'$ when the
disk has rotated through an angle $\omega t'$. Viewed from the underlying
inertial system, the signal path is  a straight line from point \textbf{2}
$(r,-\theta)$ to point \textbf{1} $(a,\omega t')$ subtending an
angle $\theta+\omega t'$ at the centre \textbf{0}. But \textbf{3}
$(r,\theta+\omega t')$ at time $t'$ and \textbf{1} $(a,0)$ at time
$t=0$ also subtend an angle $\theta+\omega t'$ at \textbf{0} in
the underlying inertial system. In the underlying inertial system
the line segments $\boldsymbol{1}_{0}\boldsymbol{3}_{t'}$ and $\boldsymbol{2}_{0}\boldsymbol{1}_{t'}$
are of equal length (indeed reflections of each other in the bisector
of the arc $\boldsymbol{1}_{0}\boldsymbol{1}_{t'}$) so that $t_{13}=t_{21}$.
Similarly, if a signal from \textbf{3} at time $t=0$ reaches \textbf{1}
at time $t=t''$ then $\boldsymbol{3}_{0}\boldsymbol{1}_{t''}$ and
$\boldsymbol{1}_{0}\boldsymbol{2}_{t''}$ are of equal length and
$t_{31}=t_{12}$. \end{shaded}%
\end{minipage}

\subsubsection{Nesting commands}

Analytically, both $t_{21}$ and $t_{13}$ are the same fixed point
of the function of $t$
\[
c^{-1}\sqrt{r^{2}+a^{2}-2ra\cos(\theta+\omega t)}
\]
 and $t_{31}$ and $t_{12}$ are the same fixed point of the function
of $t$
\[
c^{-1}\sqrt{r^{2}+a^{2}-2ra\cos(\theta-\omega t).}
\]
To calculate $2t_{12}+t_{23}$ therefore means calculating
\begin{verbatim}
    2\iter*{ c^{-1}\sqrt{a^2+r^2-2ar
                  \cos(\theta-\omega t)} }
      + \iter*{ c^{-1}\sqrt{2r^2-2r^2
                  \cos(2\theta+\omega t)} }
\end{verbatim}
with the analogous expression for $2t_{13}+t_{32}$. But we can do
the comparison of round trip times `in one go' by nesting the \verb`\iter*`
commands inside an \verb`\eval*` command:
\begin{verbatim}
     
    \eval*{ % circuit 1231
        2\iter[var=t]{ c^{-1}\sqrt{a^2+r^2-2ar
                  \cos(\theta-\omega t)} }[8]
      + \iter[var=t]{ c^{-1}\sqrt{2r^2-2r^2
                  \cos(2\theta+\omega t)} }[8]
    % circuit 1321
      - 2\iter[var=t]{ c^{-1}\sqrt{a^2+r^2-2ar
                  \cos(\theta+\omega t)} }[8]
      - \iter[var=t]{ c^{-1}\sqrt{2r^2-2r^2
                  \cos(2\theta-\omega t)} }[8]
          }[ a=10,r=20,\theta=0.2,t=1 ] 
\end{verbatim}
$\Longrightarrow$     \eval*{ % circuit 1231
        2\times\iter[var=t]{ c^{-1}\sqrt{a^2+r^2-2ar
                  \cos(\theta-\omega t)} }[8]
      + \iter[var=t]{ c^{-1}\sqrt{2r^2-2r^2
                  \cos(2\theta+\omega t)} }[8]
    % circuit 1321
      - 2\times\iter[var=t]{ c^{-1}\sqrt{a^2+r^2-2ar
                  \cos(\theta+\omega t)} }[8]
      - \iter[var=t]{ c^{-1}\sqrt{2r^2-2r^2
                  \cos(2\theta-\omega t)} }[8]
          }[a=10,r=20,\theta=0.2,{t}=1] . 

By itself this result is of little interest beyond seeing that \texttt{numerica-plus}
can handle the calculation. What \emph{is} interesting is to find
values of our parameters for which the time difference vanishes \textendash{}
say values of $\theta$, given the other parameters and especially
the value of $r$. Is there a circuit such that it takes a signal
the same time to travel in opposite senses around the circuit, despite
the rotation of the disk? Rather than nesting the \verb`\iter` commands
inside an \verb`\eval`, we need to nest them in a \verb`\solve`
command:
\begin{verbatim}
    \solve[env=multline*,p=.,var=\theta,+=1]
      {% circuit 1231    
        2\times\iter[var=t,+=1]{ c^{-1}\sqrt{a^2+r^2-2ar
                  \cos(\theta-\omega t)} }
      + \iter[var=t,+=1]{ c^{-1}\sqrt{2r^2-2r^2
                  \cos(2\theta+\omega t)} }
      % circuit 1321
      - 2\times\iter[var=t,+=1]{ c^{-1}\sqrt{a^2+r^2-2ar
                  \cos(\theta+\omega t)} }
      - \iter[var=t,+=1]{ c^{-1}\sqrt{2r^2-2r^2
                  \cos(2\theta-\omega t)} }
      }[ a=10,r=20,\theta=0.1,t=1 ][4]
\end{verbatim}
$\Longrightarrow$     \solve[env=multline*,p=.,var=\theta,+=1]
      {% circuit 1231    
        2\times\iter[var=t,+=1]{ c^{-1}\sqrt{a^2+r^2-2ar
                  \cos(\theta-\omega t)} }
      + \iter[var=t,+=1]{ c^{-1}\sqrt{2r^2-2r^2
                  \cos(2\theta+\omega t)} }
      % circuit 1321
      - 2\times\iter[var=t,+=1]{ c^{-1}\sqrt{a^2+r^2-2ar
                  \cos(\theta+\omega t)} }
      - \iter[var=t,+=1]{ c^{-1}\sqrt{2r^2-2r^2
                  \cos(2\theta-\omega t)} }
      }[ a=10,r=20,\theta=0.1,t=1 ][4]

One point to note here is the use of \verb`\times` (in \verb`2\times\iter`).
In this example the formula is displayed \textendash{} because of
the use of the \verb`env` setting, \verb`env=multline*`. Without
the \verb`\times` the result would have been the same but the display
of the formula would have juxtaposed the `$2$'s against the following
decimals, making it look as if signal travel times were $20.5378$
and $20.6144$ (and no doubt causing perplexity). The unfamiliar settings
are discussed in the relevant chapters below.

So this expression gives a value of $\theta_{\Delta t=0}$ for one
value of $r$. The obvious next step is to create a table of such
values, which can be done with the \verb`\tabulate` command from
the associated package \texttt{numerica-tables} wrapped around an
expression like the one above. (For my own interest I have done this.
On a High St laptop it is not fast \textendash{} plenty of time to
make a nice hot cup of tea.) But this is not a research paper on the
rotating disk. I wished to show how the different commands of \texttt{numerica-plus}
can be used to explore a meaningful problem. And although it looks
as if a lot of typing is involved, once $c^{-1}\sqrt{r^{2}+a^{2}-2ra\cos(\theta-\omega t)}$
has been formed in \LaTeX{} and values specified in the vv-list (and
the \verb`\constants` command in the preamble), much of the rest
is copy-and-paste with minor editing.

\section{Shared syntax of the new commands}

\texttt{numerica-plus} offers three new commands for three processes:
\verb`\nmcIterate` (short-name form \verb`\iter`) for iterating
functions, \verb`\nmcSolve` (short-name form \verb`\solve`) for
finding the zeros or (local) extrema of functions, and \verb`\nmcRecur`
(short-name form \verb`\recur`) for calculating terms of recurrence
relations. All three commands share the syntax of \verb`\nmcEvaluate`
(or \verb`\eval`) detailed in the associated document \texttt{numerica.pdf}.
When all options are used the command looks like, for instance,
\begin{centred}
\noindent\verb`\nmcIterate*[settings]{expr.}[vv-list][num. format]`
\end{centred}
You can substitute \verb`\nmcSolve`, or \verb`\nmcRecur` for \verb`\nmcIterate`
here. The arguments are the same as those for \verb`\nmcEvaluate`. 
\begin{enumerate}
\item \verb`*` optional switch; if present ensures a single number output
with no formatting, or an appropriate error message if the single
number cannot be produced;
\item \verb`[settings]` optional comma-separated list of \emph{key=value
}settings for this particular command and calculation;
\item \verb`{expr.}` the only mandatory argument; the mathematical expression
in \LaTeX{} form that is the object of interest;
\item \verb`[vv-list]` optional comma-separated list of \emph{variable=value
}items; for \verb`\iter` and \verb`\solve` the \emph{rightmost}
(or innermost) variable in the vv-list may have special significance;
\item \verb`[num. format]` optional format specification for presentation
of the numerical result (rounding, padding with zeros, scientific
notation); boolean output is suppressed for these commands.
\end{enumerate}
The way the result is displayed follows the same pattern as for \verb`\nmcEvaluate`
(see the associated document \texttt{numerica.pdf}), depending on
whether a math environment wraps around a command, is wrapped within
a command, is invoked with the \verb`env` setting, or is completely
absent. And with version $3$ of \texttt{numerica-plus}, as in \texttt{numerica},
there is always the \verb`f` setting which turns on (\verb`f=1`)
and turns off (\verb`f=0`) display of the formula. These matters
are irrelevant for the starred forms of commands, which give number-only
results. Looking at the various examples in the preceding section
on the rotating disk you will see illustrations of many of these different
situations. 

As well as the \verb`f` setting, most of the settings available to
the \verb`\eval` command are also available to the present commands
although not all will be relevant or have effect.\footnote{In particular the \texttt{ff} (multi-formula) setting has not, as
yet, been implemented in version 3.0.0 of \texttt{numerica-plus} (in
the interest of getting the whole revised \texttt{numerica} suite
launched).} Refer to the associated document \texttt{numerica.pdf} for a discussion
of such settings. The `trick' in version 2 of \texttt{numerica}
and \texttt{numerica-plus} of converting environments delimited by
\verb`\[` and \verb`\]`) into \verb`multline` environments whenever
the vv-list specification included a newline character (\verb`\\`)
has been dispensed with. Now, if you want a \verb`multline` or other
environment, directly specify it with the \verb`env` setting. The
enhanced handling of environments in \texttt{numerica} version 3 means
the \verb`*` setting which was available in earlier versions of \texttt{numerica}
and \texttt{numerica-plus} to suppress equation numbering is now unnecessary
and has been removed (although its presence will not cause an error
\textendash{} only leave a message in the log file and on the terminal).
Starring the environment in the standard \LaTeX{} way, \verb`env=equation*`,
\verb`env=multline*`, etc., does the job. 

In addition to the inherited settings, each of the \texttt{numerica-plus}
commands has settings of its own, discussed at the relevant parts
of the following chapters.

Section \ref{sec:introExampleOfUse} also provides examples of commands
being nested. Nesting may occur not only in the main argument of a
command, but also in the vv-list, or even the settings option. (The
associated document \texttt{numerica.pdf} has an example of the last
possibility.)

\chapter{Iterating functions: \texttt{\textbackslash nmcIterate}}

\label{chap:iterIterating-functions}Only in desperation would one
try to evaluate a continued fraction by stacking fraction upon fraction
upon fraction like so:
\begin{verbatim}
    \eval{\[ 1+\frac{1}{1+\frac{1}{1+\frac{1}
      {1+\frac{1}{1+\frac{1}{1+\frac{1}
        {1+\frac{1}{1+\frac{1}{1+\frac{1}
          {1+\frac{1}{1+\frac{1}
            {1+\frac{1}{1}}}}}}}}}}}} \]}
\end{verbatim}
$\Longrightarrow$     \eval{\[ 1+\frac{1}{1+\frac{1}{1+\frac{1}
      {1+\frac{1}{1+\frac{1}{1+\frac{1}
        {1+\frac{1}{1+\frac{1}{1+\frac{1}
          {1+\frac{1}{1+\frac{1}
            {1+\frac{1}{1}}}}}}}}}}}} \]}

\noindent\texttt{numerica-plus} provides a command for tackling problems
like this sensibly. In such problems a function is repeatedly applied
to itself (iterated). This is done through the command \verb`\nmcIterate`
or (short-name form) \verb`\iter`. 

\section{Basic use}

\noindent\label{sec:iterBasic-use}To evaluate the continued fraction
we want to apply $1+1/x$ to itself repeatedly. So, we wrap \verb`\iter`
around \verb`1+1/x` and give \verb`x` an initial value \verb`1`
in the vv-list: 
\begin{centred}
\verb`\iter{\[ 1+1/x \]}[x=1]` $\Longrightarrow$ \iter{\[ 1+1/x \]}[x=1]
\end{centred}
This hints that it might be heading in the direction of \eval{$ \tfrac{\sqrt{5}+1}2 $}
but clearly not enough iterations have been performed to confirm this.
By default, the \verb`\iter` command performs $5$ evaluations, the
initial evaluation producing in this instance the result $2$, which
then becomes the new value of $x$, producing the first iteration
proper and giving the value $1.5$, which in turn becomes the new
value of $x$ and so on. Also by default, \verb`\iter` displays the
first evaluation and the results of the final $4$. 

Both these numbers, $5$ and $4$, are likely to be too small. They
can be easily changed with the \verb`do` and \verb`see` settings.
Increasing the number of iterations in the example to \verb`do=17`
and the displayed results to \verb`see=5` shows how the iteration
of $1+1/x$ indeed stabilizes at $1.618034$:
\begin{centred}
\verb`\iter[do=17,see=5]{\[ 1+1/x \]}[x=1]` $\Longrightarrow$ \iter[do=17,see=5]{\[ 1+1/x \]}[x=1] 
\end{centred}
But iteration of functions is not limited to continued fractions.
Particularly since the emergence of chaos theory, iteration has become
an important study in its own right. Any function with range within
its domain can be iterated \textendash{} repeatedly applied to itself.
The cosine is an example:
\begin{centred}
\verb`\iter[do=20]{\[ \cos x \]}[x=\pi/2]` $\Longrightarrow$ \iter[do=20,see=4]{\[ \cos x \]}[x=\pi/2]
\end{centred}
which displays the initial value and last four of 20 evaluations of
$\cos x$ when the initial value of $x$ is $\tfrac{\pi}{2}$. It
looks as if the cosine is `cautiously' approaching a limit, perhaps
around $0.738$ or $0.739$. We need to nearly double the number of
iterations to confirm that this is so (to the default $6$ figures):
\begin{centred}
\verb`\iter[do=39]{\[ \cos x \]}[x=\pi/2]` $\Longrightarrow$ \iter[do=39]{\[ \cos x \]}[x=\pi/2]
\end{centred}
The display is largely fixed, hard-coded, at least at this stage.
It uses an \verb`array` environment. Whether it is centred or treated
in an inline manner depends on whether displaystyle or inline (or
no) delimiters are used. Rather than using explicit delimiters, the
\verb`env` setting can also be used (there are examples below).

As already noted, for a function to be iterated indefinitely, its
range must lie within or be equal to its domain. If even part of the
range of a function lies outside its domain then on repeated iteration
there is a chance that a value will eventually be calculated which
lies in this `outside' region. Iteration cannot continue beyond
this point and an error message is generated. As an example consider
the inverse cosine, \verb`\arccos`. This can be iterated only so
far as the iterated values lie between $\pm1$ inclusive. If we try
to iterate \verb`\arccos` at 0 for example, since $\cos\frac{1}{2}\pi=0$,
$\arccos0=\eval{0.5\pi}[4]$ ($\tfrac{1}{2}\pi$) so only a first
iterate is possible. But we could choose an initial value more carefully;
$37$ iterations of the cosine at $\tfrac{1}{2}\pi$ led to a fixed
point $0.739085$, so let's choose $0.739085$ as initial point and
perform $37$ iterations:
\begin{centred}
\verb`\iter[do=37]{\[ \arccos x \]}[x=0.739085]` $\Longrightarrow$
\iter[do=37]{\[ \arccos x \]}[x=0.739085]
\end{centred}
The result of the $37$th iteration is greater than $1$. Thus increasing
the number of iterations to 38 should generate an error message:
\begin{centred}
\verb`\iter[do=38,see=4]{\[ \arccos x \]}[x=0.739085]` $\Longrightarrow$\iter[do=38,see=4]{\[ \arccos x \]} [x=0.739085]  
\end{centred}
which it does. \verb`l3fp` objects when asked to find the inverse
cosine of a number greater than $1$.

\subsection{Logistic map}

\label{subsec:iterLogistic-map}The logistic map came to prominence
with a 1976 paper by the biologist Robert May. He examined the equation
\[
x_{n+1}=rx_{n}(1-x_{n}),
\]
where $x_{n}\in[0,1]$ and represents the ratio of an existing population
to the maximum possible population. The intent is to capture two opposed
effects: reproduction, with the rate of population increase proportional
to the population when the population is small ($x_{n+1}\approx rx_{n}$),
and mortality, when the population approaches the `carrying capacity'
of the environment ($x_{n+1}\approx r(1-x_{n})$ ).

The logistic map $rx(1-x)$ exhibits a variety of behaviours depending
on the value of $r\in(0,4)$. The Wikipedia article on the subject
lists the following:
\begin{enumerate}
\item With $0\le r\le1$, the population dies, independent of the initial
population.
\item With $1\le r\le2$, the population will quickly approach the value
$1-1/r$, independent of the initial population.
\item With $2\le r\le3$, the population will also eventually approach the
same value $1-1/r$, but first will fluctuate around that value for
some time. The rate of convergence is linear, except for $r=3$, when
it is dramatically slow.
\item With $r$ between $3$ and $1+\surd6\approx3.44949$ the population
will eventually oscillate between two values $x_{\pm}=\bigl(r+1\pm\sqrt{(r-3)(r+1)}\,\bigr)\big/2r$.
\item With $r$ between $3.44949$ and $3.54409$ (approximately), from
almost all initial conditions the population eventually oscillates
among four values.
\item With $r$ beyond $3.54409$, from almost all initial conditions, the
population eventually oscillates among $8$ values, then $16$, then
$32$, and so on. The lengths of the intervals of $r$ for each oscillation
regime decrease rapidly; the ratio between the lengths of successive
bifurcation intervals approaches the Feigenbaum constant $\delta\approx4.66920$.
This is an example of a period-doubling cascade.
\item $r\approx3.56995$, at the end of the period-doubling cascade, marks
the onset of chaos. Most values of $r$ beyond this, from almost all
initial conditions, no longer yield periodic oscillations. Slight
variations in the initial population value yield significantly different
results over time, but there are still certain isolated ranges of
$r$ (islands of stability) that show non-chaotic behavior.
\end{enumerate}
This is a rich landscape of possibilities to explore. For instance,
with $r=3.2$ we get a period-$2$ cycle, oscillating between $x_{\pm}$
with values
\begin{verbatim}
    \eval[sep=\ \text{and}\ ,p=.,ff]
        { (1/2r)[r+1+\sqrt{(r-3)(r+1)}] , 
          (1/2r)[r+1-\sqrt{(r-3)(r+1)}] }[r=3.2]
\end{verbatim}
$\Longrightarrow$ \eval[sep=\ \text{and}\ ,p=.,ff]
        { (1/2r)[r+1+\sqrt{(r-3)(r+1)}] , 
          (1/2r)[r+1-\sqrt{(r-3)(r+1)}] }[r=3.2] 
\begin{centred}
\verb`\iter[do=18,see=6]{\[ rx(1-x) \]}[r=3.2,x=0.5]` $\Longrightarrow$
\iter[do=18,see=6]{\[ rx(1-x) \]}[r=3.2,x=0.5]
\end{centred}
With $r=3.6$ we get chaos. For initial value of $x$ I have chosen
neighbouring values, $0.3$ and $0.31$:
\begin{centred}
\verb`\iter[env=\[,do=100,see=6]{ rx(1-x) }[r=3.6,x=0.3]` $\Longrightarrow$
 \iter[env=\[,do=100,see=6]{ rx(1-x) }[r=3.6,x=0.3]

\verb`\iter[env=\[,do=100,see=6]{ rx(1-x) }[r=3.6,x=0.31]` $\Longrightarrow$
 \iter[env=\[,do=100,see=6]{ rx(1-x) }[r=3.6,x=0.31]
\end{centred}
If you are using the decimal comma, make sure that the variables in
the vv-list are separated by \emph{semicolons} rather than commas,
otherwise puzzling errors are almost certain to arise.

\section{Star (\texttt{{*}}) option: fixed points}

In some of the preceding examples, iteration eventually ended at a
\emph{fixed point} \textendash{} a point $x$ where $f(x)=x$. Appending
a star (asterisk) to the \verb`\iter` command is the signal for the
\verb`\iter` command to continue iterating until a fixed point is
reached at the specified rounding value, or some fixed maximum number
of iterations have been performed. The star overrides any value specified
by the the \verb`do` setting. It also overrides any elements of the
display other than the numerical result \textendash{} meaning negative
results display with a hyphen for the minus sign unless \verb`\iter*`
is placed in a math environment. 

Return to an earlier example, the continued fraction $1+1/x$. Starring
the \verb`\iter` command gives 
\begin{centred}
\verb`\iter*{ 1+1/x }[x=1]` $\Longrightarrow$ \iter*{ 1+1/x }[x=1],
\end{centred}
and generalizing the example,
\begin{centred}
\verb`\iter*{ 1+a/x }[a=n(n+1),n=3,x=1]` $\Longrightarrow$ \iter*{ 1+a/x }[a=n(n+1),n=3,x=1]
\end{centred}
Indeed, trying in turn $n=0,1,2,3,4,5$ we see that when iterated
$1+a/x\to n+1$. A little invesigating shows that this is hardly surprising.
If $x$ is a point such that 
\[
1+\frac{n(n+1)}{x}=x
\]
then $x^{2}-x-n(n+1)=0$. The quadratic factorizes: $(x-(n+1))(x+n)=0$
so that, indeed, $x=n+1$ is a fixed point, as also \textendash{}
we learn \textendash{} is $x=-n$, trivially. There is nothing here
that requires $n$ to be an integer,
\begin{centred}
\verb`\iter*{ 1+a/x }[a=n(n+1),n=(\surd5-1)/2,x=1]` $\Longrightarrow$
\iter*{ 1+a/x }[a=n(n+1),n=(\surd5-1)/2,x=1],
\end{centred}
but if we put $n=6$, we get a message:
\begin{centred}
\verb`\iter*{ 1+a/x }[a=n(n+1),n=6,x=1]` $\Longrightarrow$ \iter*{ 1+a/x }[a=n(n+1),n=6,x=1]
\end{centred}
It is easy to increase the $100$ here to a larger value as we will
do in a moment but it is worth using the \verb`\info` command from
\texttt{numerica} to see what is happening. For $n=1$ and initial
value $x=1$,
\begin{verbatim}
    \iter*{1+a/x}[a=n(n+1),n=1,x=1],
    we see that it takes \info{iter}
\end{verbatim}
$\Longrightarrow$ \iter*{1+a/x}[a=n(n+1),n=1,x=1],
we see that it takes \info{iter} to attain a $6$-figure fixed point; when $n=2$ for this same initial
value it takes 41; \ldots{} ; and when $n=5$ it takes $96$. The message
when $n=6$ is hardly surprising.

The maximum prevents \verb`\iter` falling into an infinite loop or
similar state. We saw with the logistic map that there are parameter
values that lead to $2$-cycles, $4$-cycles, $8$-cycles \ldots{}
chaos, where iteration would continue indefinitely if there were no
safeguard like a specified maximum number of iterations. For our case,
however, it doesn't look as if infinite loops of this kind are the
problem. We increase the maximum by using the setting \verb`max`:
\begin{centred}
\verb`\iter*[max=150]{1+a/x}[a=n(n+1),n=6,x=1], taking \info{iter}`
$\Longrightarrow$ \iter*[max=150]{1+a/x}[a=n(n+1),n=6,x=1], taking \info{iter},
\end{centred}
and the fixed point is safely attained well within the bound of the
new maximum. 

Alternatively, we could have reduced the rounding value, say from
the default $6$ to $5$:
\begin{centred}
\verb`\iter*{1+a/x}[a=n(n+1),n=6,x=1][5], taking \info{iter}` $\Longrightarrow$
\iter*{1+a/x}[a=n(n+1),n=6,x=1][5], taking \info{iter}.
\end{centred}
A fixed point is attained \textendash{} but with no room to spare.
Generally, reducing the rounding value is the other strategy to pursue
when faced with the `No fixed point attained' message, and perhaps
the better one initially. If there is no fixed point after $100$
iterations at some low rounding value \textendash{} say $2$ or $3$
\textendash{} then there may well be no fixed point at all. 

\verb`\iter` determines that a fixed point has been attained when
the difference between successive iterations vanishes when rounded
to the current rounding value. This can lead to an error in the final
digit: the \emph{difference }may vanish but the final value round
\emph{away} from the fixed point. To seek reassurance that the fixed
point really is the correct value you might wish to seek a fixed point
at a higher rounding value without changing the number of digits displayed.
The extra rounding is achieved by entering \verb`+=<integer>` in
the settings option, where \verb`<integer>` is the \emph{extra} rounding
desired.

For example, for the logistics map with $r=1.5$ we expect a fixed
point with value 
\begin{centred}
\verb`\eval{1-1/r}[r=1.5]` $\Longrightarrow$ \eval{1-1/r}[r=1.5],
\end{centred}
and indeed
\begin{centred}
\verb`\iter*{rx(1-x)}[r=1.5,x=0.5]` $\Longrightarrow$ \iter*{rx(1-x)} [r=1.5,x=0.5],
\end{centred}
the expected value \textendash{} if we ignore the final digit. So
let's use the \verb`+` setting, to demand that the difference between
successive iterate values at $6+1=7$ figures vanishes. That should
ensure the correct $6$-figure fixed point is attained, and it is:
\begin{centred}
\verb`\iter*[+=1]{rx(1-x)}[r=1.5,x=0.5]` $\Longrightarrow$ \iter*[+=1]{rx(1-x)}[r=1.5,x=0.5].
\end{centred}

\subsection{Use with \texttt{\textbackslash nmcInfo}}

We have already seen that the \verb`\nmcInfo` command provides information
on the number of iterations necessary to attain a fixed point, especially
how many are required at a particular rounding value. That knowledge
allows a good guess as to whether a fixed point will be attained at
a greater rounding value. Thus when iterating the function
\[
f(t_{ij})=c^{-1}\sqrt{r_{i}^{2}+r_{j}^{2}-2r_{i}r_{j}\cos(\theta_{j}-\theta_{i}+\omega t_{ij})}
\]
in �\ref{sec:introExampleOfUse} only $5$ iterations were required
to attain $6$-figure accuracy for the fixed point. That information
came by following the \verb`\iter*` command with \verb`\nmcInfo`
(or \verb`\info`) with the argument \verb`iter`. And generally,
for any `infinite' process, follow the command with an \verb`\info`
command if you want to know how many `steps' \textendash{} in the
present case iterations \textendash{} are required to achieve the
result. So, if $5$ iterations achieve $6$-figure accuracy, presumably
something like $10$ iterations will achieve $12$-figure accuracy:
\begin{verbatim}
    \iter*{ c^{-1}\sqrt{r_i^2+r_j^2-2r_i r_j
      \cos(\theta_{ij}+\omega t)}
        }[ r_i=10,r_j=20,\theta_{ij}=0.2,t=1 ][12]
    ,\quad\info{iter}. 
\end{verbatim}
$\Longrightarrow$     \iter*{ c^{-1}\sqrt{r_i^2+r_j^2-2r_i r_j
      \cos(\theta_{ij}+\omega t)}
        }[ r_i=10,r_j=20,\theta_{ij}=0.2,t=1 ][12]         
    ,\quad\info{iter}. (Remember, \texttt{numerica-plus} knows the values of $c$ and $\omega$
from a \verb`\constants` statement in the preamble.) And indeed only
$9$ iterations suffice to achieve $12$-figure accuracy: 
\begin{verbatim}
    \iter[env=\[,vv=,do =11,see=4]
      { c^{-1}\sqrt{r_i^2+r_j^2-2r_i r_j
        \cos(\theta_{ij}+\omega t)}
      }[ r_i=10,r_j=20,\theta_{ij}=0.2,t=1 ][12]
\end{verbatim}
$\Longrightarrow$     \iter[env=\[,vv=,do =11]
      { c^{-1}\sqrt{r_i^2+r_j^2-2r_i r_j
        \cos(\theta_{ij}+\omega t)}
      }[ r_i=10,r_j=20,\theta_{ij}=0.2,t=1 ][12]

\noindent (Display of the vv-list has been suppressed with the setting
\verb`vv=` .)

Or again, with another example from earlier,
\begin{centred}
\verb`\iter*{\cos x}[x=\pi/2],\ \info{iter}` $\Longrightarrow$ \iter*{\cos x}[x=\pi/2],\ \info{iter}.
\end{centred}
That suggests that around $2\times37=74$ iterations will give a $2\times6=12$-figure
answer, well within the cut-off figure of $100$:
\begin{centred}
\verb`\iter*{\cos x}[x=\pi/2][12],\ \info{iter}.` $\Longrightarrow$
\iter*{\cos x}[x=\pi/2][12],\ \info{iter}.
\end{centred}

\section{Settings, saving results, errors, }

\label{sec:iterSettings-option}The settings option is a comma-separated
list of items of the form \emph{key~=~value}. Only some of the keys
for \verb`\nmcEvaluate` discussed in Chapter~$5$ of the associated
document \texttt{numerica.pdf} are relevant for \verb`\nmcIterate`.
Thus should a quantity in the vv-list depend on the iteration variable,
forcing an implicit mode calculation, simply enter \verb`vv@=1` (alternatively,
\verb`vvmode=1`) in the settings option, as with \verb`\eval`:
\begin{centred}
\verb`\iter*[vv@=1]{ f(x) }[f(x)=1+a/x,a=12,x=1]` $\Longrightarrow$
\iter*[vv@=1]{ f(x) }[f(x)=1+a/x,a=12,x=1].
\end{centred}
Implicit in the example is the default multi-token setting \verb`xx=1`
inherited from \verb`\eval` and ensuring that the multi-token variable
$f(x)$ is treated correctly. 

We could add \verb`dbg=1` to the example \textendash{} or just enter
\verb`view` \textendash{} to get a glimpse at the `innards' of
what is going on:
\begin{verbatim}
    \iter*[view,vv@=1]{ f(x) }[f(x)=1+a/x,a=12,x=1]
    \info{iter}
\end{verbatim}
$\Longrightarrow$ \iter*[view,vv@=1]{ f(x) }[f(x)=1+a/x,a=12,x=1]
    \info{iter}

The multi-token variable \verb`f(x)` has been changed to a single-token.
The difference between the two long `stored' numbers is less than
$5$ in the $7$th decimal place, meaning a fixed point has been found.
Since $59$ iterations are required to attain the fixed point at $6$-decimal
place accuracy, the values shown correspond to the $60$th iteration,
the \emph{final} iteration. Because the \verb`\iter` command is the
starred form, the result that is fed to \LaTeX{} is simply the fixed
point $4$ expressed as a number. Remove the star, add \verb`do=60`
and replace \verb`view` with \verb`dbg=55` (or equivalently \verb`dbg=5*11`)
in the settings option:
\begin{centred}
\verb`\iter[dbg=55,vv@=1,do=60]{ f(x) }[f(x)=1+a/x,a=12,x=1]` $\Longrightarrow$
\iter[dbg=55,vv@=1,do=60]{ f(x) }[f(x)=1+a/x,a=12,x=1]
\end{centred}
Now the \LaTeX{} form is much fuller and the stored numbers exactly
match those of the starred form.

\subsection{\texttt{\textbackslash nmcIter}ate-specific settings}

\begin{table}
\centering{}\caption{\protect\label{tab:iterSettings}Settings for \texttt{\textbackslash nmcIterate}}
\begin{center}
\begin{tabular}{llll}
\toprule 
{\small key} & {\small type} & {\small meaning} & {\small initial}\tabularnewline
\midrule
{\small\texttt{var}} & {\small token(s)} & {\small iteration variable} & \tabularnewline
{\small\texttt{+}} & {\small int} & {\small fixed point extra rounding} & {\small\texttt{0}}\tabularnewline
{\small\texttt{max}} & {\small int > 0} & {\small max. iteration count (fixed points)} & {\small\texttt{100}}\tabularnewline
{\small\texttt{do}} & {\small int > 0} & {\small number of iterations to perform} & {\small\texttt{5}}\tabularnewline
{\small\texttt{see}} & {\small int > 0} & {\small number of final iterations to view} & {\small\texttt{4}}\tabularnewline
 & {\small int ($\mathtt{0}/\mathtt{1}/\mathtt{2}$)} & {\small form of result saved with }{\small{\small\verb`\`}} & {\small\texttt{0}}\tabularnewline
\bottomrule
\end{tabular}
\par\end{center}
\end{table}
In addition to the inherited settings there are some specific to \verb`\nmcIterate`.
These are listed in Table~\ref{tab:iterSettings}.

\subsubsection{Iteration variable}

In nearly all of the examples so far, the iteration variable has been
the rightmost variable in the vv-list and has not needed to be otherwise
specified. However it is sometimes not feasible to indicate the variable
in this way. In that case, entering
\begin{verbatim}
    var = <variable name>
\end{verbatim}
in the settings option enables the variable to be specified, irrespective
of what the rightmost variable in the vv-list is. Here, \verb`<variable name>`
will generally be a character like \verb`x` or \verb`t` or a token
like \verb`\alpha`, but it could also be a multi-token name like
{\ttfamily\verb`x'`}\texttt{ }or \verb`\beta_{ij}` (or
even \verb`Fred` if you so chose). Although the iteration variable
can be independently specified like this, it must still be given an
initial \emph{value} in the vv-list \textendash{} only now it need
not be the rightmost variable.

In the following example the rightmost variable is \verb`n` which
is clearly \emph{not} the iteration variable:
\begin{centred}
\verb`\iter[var=x,do=40]{$ 1+a/x $}[x=n-1,a=n(n+1),n=2][*]` $\Longrightarrow$
\iter[var=x,do=40]{$ 1+a/x $}[x=n-1,a=n(n+1),n=2][*]
\end{centred}

\subsubsection{Extra rounding for fixed-point calculations}

\label{subsec:iterExtra-rounding}The criterion used to signal the
attainment of a fixed point is that the difference between successive
iterations vanishes when rounded to the current rounding value. As
already noted, because of rounding errors and the `round to even'
tie-breaking rule, this criterion may lead to error in the last digit.
That can be avoided by rounding to a greater number of digits than
the actual rounding value. This \emph{extra} rounding is achieved
by entering 
\begin{verbatim}
    + = <integer>
\end{verbatim}
in the settings option. By default this extra rounding is set to zero. 

We have seen before that $\cos x$ starting at $x=\tfrac{1}{2}\pi$
takes $37$ iterations to reach a $6$-figure fixed point $0.739085$,
about $6$ iterations per decimal place. By entering \verb`+=1` in
the settings option the number of iterations is increased to $43$,
$6$ more than $37$ but, reassuringly, the $6$-figure result that
is displayed remains unchanged:
\begin{centred}
\verb`$ \iter*[+=1]{\cos x}[x=\pi/2] $,\ \info{iter}` $\Longrightarrow$
$ \iter*[+=1]{\cos x}[x=\pi/2] $,\ \info{iter}.
\end{centred}

\subsubsection{Maximum {\small iteration count for fixed-point searches}}

\label{subsec:iterMaximum-iteration-count}To prevent a fixed-point
search from continuing indefinitely, perhaps because no fixed point
exists, there needs to be a maximum number of iterations specified
after which point the search is called off. By default this number
is $100$. To change it enter
\begin{verbatim}
    max = <positive integer>
\end{verbatim}
in the settings option.

\subsubsection{Number of iterations to perform}

To specify the number of iterations to perform enter
\begin{verbatim}
    do = <positive integer>
\end{verbatim}
in the settings option. Note that if the \verb`*` option is present
this value will be ignored and iteration will continue until either
a fixed point or the maximum iteration count is reached. By default
\verb`do` is set to $5$. (Note that \texttt{do} can be set to a
greater number than the initial $100$ setting of \verb`max`; \verb`max`
applies only to the starred form \verb`\iter*`.)

\subsubsection{Number of iterations to show}

To specify the number of final iterations to show enter 
\begin{verbatim}
    see = <positive integer>
\end{verbatim}
in the settings option. By default \verb`see` is set to $4$. Always
it is the \emph{last} \verb`see` iterations that are displayed. If
\verb`see` is set to an equal or greater value than \verb`do`, all
iterations are shown. If the star option is used the \verb`see` value
is ignored.

\subsection{Form of result saved by \texttt{\textbackslash nmcReuse}}

In version 2 of \texttt{numerica-plus} there was a setting \verb`reuse`
that determined the content of what was saved with the\texttt{ }\verb`\nmcReuse`
command. This has been removed. Now it is only the last \verb`see`
values that are saved as a comma list. The values stored in the list
are each wrapped in braces. In this way no confusion arises if the
decimal comma is being used. For a fixed point calculation, it is
the fixed point as displayed that is saved. 
\begin{verbatim}
    \iter[do=12,see=4]
      {\[ kx(1-x) \]}[k=3.5,x=0.5]
    \reuse{logistic}
\end{verbatim}
$\Longrightarrow$ \iter[do=12,see=4]{\[ kx(1-x) \]}[k=3.5,x=0.5] \reuse[renew]{logistic}

\noindent whence \verb`$\logistic$` $\Longrightarrow$ $\logistic$.
(I have placed \verb`\logistic` between \verb`$` delimiters to get
thin spaces after the commas.)

If you want to isolate just one member of the comma list, \texttt{numerica-plus}
offers the utility function \verb`\clitem` (`comma list item')
which acts on the control sequence followed by a number:
\begin{centred}
\verb`\clitem\logistic 3` $\Longrightarrow$ \clitem\logistic 3.
\end{centred}
Use positive numbers for counting from the left, negative numbers
for counting from the right. Multi-digit numbers must be enclosed
in braces. See �\ref{subsec:recurAccessing-individual-terms} for
more on the use of \verb`\clitem`.

\subsection{Errors}

There is only one error specifically related to \verb`\nmcIterate`
and that is when the number of iterations exceeds the specified maximum
number in a fixed point calculation. We have already met this in the
earlier discussion of fixed points. Another example is
\begin{centred}
\verb`\iter*{kx(1-x)}[k=3.5,x=0.5]` $\Longrightarrow$ \iter*{kx(1-x)}[k=3.5,x=0.5]
\end{centred}
Other \texttt{numerica} errors can also afflict an iteration calculation.
Again an example of this was provided in �\ref{sec:iterBasic-use}
when we sought to iterate \verb`\arccos` $38$ times rather than
the previously successful $37$: 
\begin{centred}
\verb`\iter[do=38,see=4]{\[ \arccos x \]}[x=0.739085]` $\Longrightarrow$\iter[do=38,see=4]{\[ \arccos x \]} [x=0.739085]  
\end{centred}
The $38$th iterate is attempting to take the inverse cosine of a
number greater than $1$; \texttt{numerica} objects.

\section{Newton-Raphson method}

\label{subsec:iterNewton-Raphson-method}This is a method for solving
equations in one variable, say $f(x)=0$, by iterating the expression
\[
x_{n+1}-\frac{f(x_{n})}{f'(x_{n})}
\]
where $f'$ is the derivative of $f$. \texttt{numerica-plus} does
not automatically calculate the derivative; given the function $f$
the user needs to manually insert (the analytic expression for) it.

\begin{wrapfigure}{o}{0.4\columnwidth}%
\centering{}\includegraphics[scale=0.8]{figures/NewtonRaphson}\caption{N-R method\protect\label{fig:Newton-Raphson-method}}
\end{wrapfigure}%
Fig.~\ref{fig:Newton-Raphson-method} depicts a situation where the
method works well. This will not always be the case. If $x=z$ is
the zero ($f(z)=0$) then should the derivative vanish at $z$ ($f'(z)=0$)
there may well be problems. I give an example below, after first giving
a simple illustration of the method.

Suppose $f(x)=\sin x$; then $f'(x)=\cos x$. If the initial test
value is $4$, then $x-\sin x/\cos x=x-\tan x$ is the expression
to insert in the \verb`\iter*` command:
\begin{centred}
\verb`\iter*{x-\tan x}[x=4], \info{iter}` $\Longrightarrow$ \iter*{x-\tan x}[x=4], \info{iter},
\end{centred}
which we recognize as the 6-figure value of $\pi$, attained very
quickly after only 3 iterations. To check this omit the asterisk:
\begin{centred}
\verb`\iter[f=1]{x-\tan x}[x=4]` $\Longrightarrow$ \iter[f=1]{x-\tan x}[x=4]
\end{centred}
But that is $4$ iterations, not the claimed $3$, before the correct
value appears. This is an example of the difference between successive
values vanishing at the $6$-figure rounding value, but the final
digit being $1$ shy of the correct value. If we put \verb`+=1` in
the settings and round to $7$ figures, we get
\begin{centred}
\verb`\iter[f=1,+=1]{x-\tan x}[x=4][7]` $\Longrightarrow$ \iter[f=1,+=1]{x-\tan x}[x=4][7]
\end{centred}
The difference between the $3$rd and $4$th values is $0.0000003$
which rounds to $0$ at $6$ figures whereas $3.1415924$ rounds to
$3.141592$.

We could also do the calculation in implicit mode (the \verb`f=1`
in the settings means the formula is part of the display and has \emph{nothing}
to do with the \verb`f` in the main argument and vv-list):
\begin{centred}
\verb`\iter[f=1,vv@=1]{x-f(x)/f'(x)}[f(x)=\sin x,{f'(x)}=\cos x,x=4]`
$\Longrightarrow$ \iter[f=1,vv@=1]{x-f(x)/f'(x)}[f(x)=\sin x,{f'(x)}=\cos x,x=4]
\end{centred}
where display of the derivative from the vv-list has been suppressed
(by the enclosing braces).

Let's tackle something less obvious, finding the roots of $\cos x\cosh x\pm1$.
(\emph{HMF }Table~4.18 gives small tables of the first five roots
in both cases; subsequent roots are essentially $\cos^{-1}(\mp1)=\tfrac{1}{2}(2n\pm1)\pi$
since $\cosh x$ is so large.) Writing $f(x)$ for the function, in
both cases $f'(x)=\cos x\sinh x-\sin x\cosh x$. In the example, our
initial test value is \verb`x=5`. We save the result in the control
sequence \verb`\zilch` and then check that it really does contain
a zero of $f(x)$ in the final statement:
\begin{verbatim}
    \iter*[vv@=1]{ x-f(x)/f'(x) }
      [f(x)=\cos x\cosh x-1, {f'(x)}=
        \cos x\sinh x-\sin x\cosh x,x=5][15],
    \quad \info{iter}
    \reuse[renew]{zilch} \macros{ \zilch } \par
    \eval{\[\cos x\cosh x -1 \]}[x=\zilch][15]
\end{verbatim}
$\Longrightarrow$      \iter*[vv@=1]{ x-f(x)/f'(x) }
      [f(x)=\cos x\cosh x-1, {f'(x)}=
        \cos x\sinh x-\sin x\cosh x,x=5][15],
    \quad \info{iter}
    \reuse[renew]{zilch} \macros{ \zilch } \par
    \eval{\[\cos x\cosh x -1 \]}[x=\zilch][15*]

It is remarkable that only $5$ iterations are required for $15$-place
accuracy; \emph{HMF} gives only $6$-figures, $4.7300407$.

\noindent{}%
\noindent\begin{minipage}[t]{1\columnwidth}%
\begin{shaded}%
Of course, $x=0$ is a zero of $\cos x\cosh x-1$, but $f'(x)$ vanishes
at zero and the Newton-Raphson method fails there for that reason.
Indeed it's easy to check that $f'(0)=f''(0)=f'''(0)=0$ and $f^{iv}(x)=-4(f(x)+1)$
meaning $f^{iv}(0)=-4$ and on repeated differentiation the pattern
repeats. Thus the Maclaurin series for $f(x)$ near $x=0$ is
\[
f(x)=\sum_{n=1}^{\infty}\frac{(-1)^{n}4^{n}x^{4n}}{(4n)!}=-\frac{x^{4}}{6}+\frac{x^{8}}{2520}-\dots
\]
\emph{Any }value of $x$ close to $0$ is going to give a value of
$f(x)=\cos x\cosh x-1$ about $4$ decimal places closer. Indeed,
when trial values like $x=1$, $x=0.5$, $x=0.3$ and so on are used
in the iteration of $x-f(x)/f'(x)$, the iteration converges on a
series of `pseudo-zeros' of $f(x)$, all artifacts of the rapidity
with which $\cos x\cosh x$ converges to $1$ at $x=0$.\end{shaded}%
\end{minipage}

\chapter{Finding zeros and extrema: \texttt{\textbackslash nmcSolve}}

\label{chap:solveSolve}\texttt{numerica-plus} provides a command,\textbf{
}\verb`\nmcSolve` (short-name form \verb`\solve`), for finding a
zero of a function, should it have one. In the following example, 
\begin{centred}
\verb`\solve[p]{\[ e^{ax}-bx^2 \]}[a=2,b=3,{x}=0]` $\Longrightarrow$
\solve[p]{\[ e^{ax}-bx^2 \]}[a=2,b=3,{x}=0]
\end{centred}
I have sought and found a solution $x$ to the equation $e^{ax/2}-bx^{2}=0$
when $a=2$ and $b=3$, starting with a trial value $x=0$, entered
as the \emph{rightmost} variable in the vv-list (and em-braced since
I don't want this trial value displaying in the presentation of the
result). Although $x$ has been found to the default six-figure accuracy,
it is evident that the function vanishes only to five figures. Let's
check:
\begin{centred}
\verb`\eval{$ bx^2 $}[b=3,x=x=-0.390647]` $\Longrightarrow$ \eval{$ bx^2 $}[b=3,x=-0.390647], 

\verb`\eval{$ e^{ax} $}[a=2,x=-0.390647]` $\Longrightarrow$ \eval{$ e^{ax} $}[a=2,x=-0.390647]; 
\end{centred}
the values agree save in the final digit.

This discrepancy in the final decimal place or places is a general
feature of solutions found by \verb`\solve`. It is the value of $x$,
not the value of $f(x)$, that is being found (in this case) to six
figures. If the graph of a function crosses the $x$-axis steeply
then the $x$ value (the zero) may be located to a higher precision
than the function value. Conversely, if the graph of a function crosses
the $x$-axis gently (at a shallow angle) then the function value
will vanish to a greater number of decimal places than the zero (the
$x$ value) is found to.

A second example, which we can check against values tabulated in \emph{HMF},
is to find a value of $x$ that satisfies $\tan x=\lambda x$. In
other words, find a zero of $\tan x-\lambda x$. In the example $\lambda$
is negative, so a trial value for $x$ greater than $\pi/2$ seems
like a good idea. I've chosen $x=2$. 
\begin{centred}
\verb`\solve{$ \tan x - \lambda x $}[\lambda=-1/0.8,{x}=2][5]` $\Longrightarrow$
\solve{$ \tan x - \lambda x $}[\lambda=-1/0.8,{x}=2][5].
\end{centred}
Table 4.19 of \emph{HMF }lists values of $x$ against $\lambda$ and
this is the $5$-decimal place value tabulated there.

\section{Extrema}

A function may not have a zero; or, for the given initial trial value
and initial step in the search for a zero, there may be a local extremum
in the way. In that case \texttt{numerica-plus} may well locate the
local extremum (maximum or minimum but not a saddle point). For example
for the quadratic $(2x-1)^{2}+3x+1$ the \verb`\solve` command gives
the result 
\begin{centred}
\verb`\solve[vv=]{$ (2x-1)^2+3x+1 $}[x=2]` $\Longrightarrow$ \solve[vvi=]{$ (2x-1)^2+3x+1 $}[{x}=2]. 
\end{centred}
Since $(2x-1)^{2}+3x+1\ne0$ for any (real number) $x$, we deduce
that the quadratic takes a minimum value $1.9375$ at $x=0.125$ \textendash{}
easily confirmed analytically. This particular minimum is a global
minimum but in general any extremum found is only \emph{local}. The
function may well take larger or smaller values (or vanish for that
matter) further afield.

It is also worth noting in this example the \verb`vv=` in the settings
option which suppresses display of the vv-list. (The only member of
the vv-list is the trial value \verb`x=2` which we do not want to
display.)

\noindent{}%
\noindent\begin{minipage}[t]{1\columnwidth}%
\begin{shaded}%
Note that the function for which a zero is being sought is \emph{not}
equated to zero when entered in the \verb`\solve` command. It is
\verb`\solve{ f(x) }`, not \verb`\solve{ f(x)=0 }`. This is precisely
because it may be an extremum that is found rather than a zero (if
extremum or zero is found at all \textendash{} think $e^{x}$). The
display of the result makes clear which is which, equating $f(x)$
to its value, zero or extremum depending on what has been found, as
you can see in the preceding examples.\end{shaded}%
\end{minipage}

\subsection{The search strategy}

\label{subsec:solveSearch-strategy}If you have some sense of where
a function has a zero, then choose a trial value in that vicinity.
\verb`\solve` uses a bisection method to home in on the zero. It
therefore needs \emph{two} initial values. For the first it uses the
trial value you specify, call it $a$ and for the second, by default,
it uses $a+1$. (The default value $1$ for the initial step from
the trial value can be changed in the settings option with the setting
\verb`dvar`; see �\ref{sec:solveSettings-option}.) If $f(a)$ and
$f(a+1)$ have opposite signs then that is good. Bisection of the
interval $[a,a+1]$ can begin immediately in order to home in on the
precise point where $f$ vanishes. Write $b=a+1$.
\begin{itemize}
\item Let $c=\tfrac{1}{2}(a+b)$; if $f(c)=0$ the zero is found; otherwise
either $f(a),f(c)$ are of opposite signs or $f(c),f(b)$ are of opposite
signs. In the former case write $a_{1}=a,$ $b_{1}=c$; in the latter
case write $a_{1}=c$, $b_{1}=b$ and then redefine $c=\tfrac{1}{2}(a_{1}+b_{1})$.
Continue the bisection process, either until an exact zero $c$ of
$f$ is reached ($f(c)=0$) or a value $c$ is reached where the difference
between $a_{n+1}$ and $b_{n+1}$ is zero at the specified rounding
value. (But note, $f(c)$ may not vanish at that rounding value \textendash{}
the zero might be elsewhere in the interval and $f$ might cross the
axis at a steep slope.)
\end{itemize}
However $f(a)$ and $f(b)=f(a+1)$ may not have opposite signs. If
we graph the function $y=f(x)$ and suppose $f(a),f(b)$ are distinct
but of the same sign, then the line through the points $(a,f(a))$,
$(b,f(b))$ will intersect the $x$-axis to the left of $a$ or the
right of $b$ depending on its slope. We search always \emph{towards
the $x$-axis} in steps of $b-a$ ($=1$ with default values).
\begin{itemize}
\item If the line intersects the axis to the left of $a$ then $c=a-(b-a)$
and we set $a_{1}=c,b_{1}=a$; if the line intersects the axis to
the right of $b$ then $c=b+(b-a)$ and we set $b_{1}=c,a_{1}=b$.
The hope is that by always taking steps in the direction towards the
$x$-axis that eventually $f(c)$ will be found to lie on the \emph{opposite}
side of the axis from $f(a_{n})$ or $f(b_{n})$, at which point the
bisection process begins.
\item Of course this may not happen. At some point $c$ may lie to the left
of $a_{n}$ but $\left|f(c)\right|>\left|f(a_{n})\right|$, or $c$
may lie to the right of $b_{n}$ but $\left|f(c)\right|>\left|f(b_{n})\right|$.
The slope has reversed. In that case we halve the step value to $\tfrac{1}{2}(b-a)$
and try again in the same direction as before from the same point
as before ($a_{n}$ or $b_{n}$ as the case may be).
\item Should we find at some point that $f(a_{n})=f(b_{n})$ then the previous
strategy does not apply. In this case we choose $a_{n+1}$ and \textbf{$b_{n+1}$}
at the quarter and three-quarter marks between $a_{n}$ and $b_{n}$.
Either $f(a_{n+1})$ and $f(b_{n+1})$ will differ and the previous
search strategy can start again or we are on the way to finding an
extremum of $f$.
\end{itemize}
As already noted it is also possible that our function has neither
zeros nor extrema. To prevent the search continuing indefinitely,
\texttt{numerica} uses a cut-off value for the maximum number of steps
pursued \textendash{} by default set at 100.

\subsection{Elusive extrema}

\label{subsec:solveElusive-extrema}The strategy `search always towards
the $x$-axis' has a consequence: it means that a local maximum above
the $x$-axis will almost certainly not be found, since `towards
the $x$-axis' pulls the search away from the maximum. Similarly
a local minimum below the $x$-axis will also not be found since `towards
the $x$-axis' pulls the search away from the minimum. 

One way of countering this elusiveness is to add a constant value
(possibly negative) to the function whose zeros and extrema are being
sought. The zeros of the function will change but the abscissae ($x$
values) of the extrema remain unchanged. If the constant is big enough
it will push a local minimum above the axis where it can be found
or, for a negative constant, push a local maximum below the axis where
it can be found.

For example $f(x)=x^{3}-x$ has roots at $-1,0,1$, a local maximum
at $-\tfrac{1}{\surd3}$ and a local minimum at $\tfrac{1}{\surd3}$.
To locate the minimum, I have added an unnecessarily large constant
$k$ to $f(x)$ ($k=1$ would have sufficed) and a start value at
the rightmost zero. Searching `towards the $x$-axis' will take
the search towards the local minimum.
\begin{centred}
\verb`\solve{$ x^3-x+k $}[k=5,{x}=1]` $\Longrightarrow$ \solve{$ x^3-x+k $}[k=5,{x}=1].
\end{centred}
Checking, \verb`\eval{$\tfrac1{\surd 3}$}` $\Longrightarrow$ \eval{$\tfrac1{\surd 3}$}.
There is a discrepancy in the $6$th decimal place which can be eliminated
by using the extra rounding setting; see �\ref{subsec:solveExtraRounding}.
The value of the local minimum of $x^{3}-x$ is then $\eval[f=1]{4.6151-5}$.

Or, we can find the value of that local minimum value and the $x$
value where it occurs `in one go' by \emph{nesting} \verb`\solve`
within the vv-list of an \verb`\eval` command:
\begin{centred}
\verb`\eval{$ x^3-x $}[x={\solve{y^3-y+k}[k=5,y=1]}]` $\Longrightarrow$
\eval{$ x^3-x $}[x={\solve[dvar=0.25]{y^3-y+k}[k=5,y=1]}].
\end{centred}
The braces around the \verb`\solve` and its vv-list hide \emph{its}
square-brackets  from the parsing of the vv-list of the \verb`\eval`
command. 

\subsection{False extrema}

A function which `has an infinity' at a particular value can result
in a false extremum being found:
\begin{centred}
\verb`\solve{$ 1/x $}[x=-1/3]` $\Longrightarrow$ \solve{$ 1/x $}[x=-1/3].
\end{centred}
One needs to look for extrema with some awareness of how the function
behaves. `Searching blind' may lead to nonsense results. In this
particular example, changing the rounding value will show the supposed
extremum jumping from one large value to another and not settling
at any particular value.

\section{Star (\texttt{{*}}) option}

A starred form of the\textbf{ }\verb`\nmcSolve` command gives a purely
numerical result; should it be negative it displays with a hyphen
for the minus sign outside a math environment. All other elements
of the display are suppressed. When commands are nested, \texttt{numerica}
and associated packages treat an inner command as if it were the starred
form, whether the star is explicitly present or not. Returning to
a previous example,
\begin{centred}
\verb`\solve*{$ \tan x - \lambda x $}[\lambda=-1/0.8,{x}=2][5]` $\Longrightarrow$
\solve*{$ \tan x - \lambda x $}[\lambda=-1/0.8,{x}=2][5],
\end{centred}
giving the zero and nothing more.

\section{Settings option}

\label{sec:solveSettings-option}The settings option is a comma-separated
key-value list of items. The keys discussed in the settings\emph{
}option for \verb`\nmcEvaluate` discussed in the associated document
\texttt{numerica.pdf }are also available for \verb`\nmcSolve`. The
very first example in this chapter used the punctuation option \texttt{p}
(\verb`\solve[p]{\[... `) to ensure a comma after the display-style
presentation of the result. We also saw in the quadratic example illustrating
extrema the use of \verb`vv=` with no value to suppress display of
the vv-list: \verb`\solve[vv=]{$ ...`. 

Putting \verb`dbg=1`, or more simply, entering \verb`view` in the
settings option produces a familiar kind of display. Using the function
\[
ct-\sqrt{a^{2}+b^{2}-2ab\cos(\beta+\omega t)}
\]
from the rotating disk problem,
\begin{verbatim}
    \solve[view,var=t]
      {$ ct-\sqrt{a^{2}+b^{2}-2ab\cos(\beta+\omega t)} $} 
        [a=10,b=20,\beta=1,{t}=0][4]
\end{verbatim}
$\Longrightarrow$     \solve[dbg=1,var=t] 
      {$ ct-\sqrt{a^{2}+b^{2}-2ab\cos(\beta+\omega t)} 
      $}[a=10,b=20,\beta=1,{t}=0][4]

\noindent If that is too much information, understand that \verb`view`
(or \verb`dbg=1`) for \verb`\nmcSolve` is equivalent to \verb`dbg=2*3*5*7*11`.
Entering \verb`dbg=2*3*5*7` for instance will trim the \LaTeX{} expression
from the debug display, and similarly, omitting other prime factors
from \verb`dbg` will result in the corresponding element being absent
from the display. (The prime factors can be multiplied out if you
wish; see Chapter~5 of the associated document \texttt{numerica.pdf}.)

\subsubsection{Multi-line display of the result}

\label{subsec:solveMulti-line-display}By default the result is presented
on a single line (unless the star option is being used). This can
be of the form \emph{function = function value, (vv-list) $\rightarrow$
variable = result}, where \emph{function value} will either be $0$
or close to it, or an extremum of the function, and \emph{result}
will be the value of the \emph{variable} producing that zero or extremum
when substituted into the function. It takes only a slightly complicated
formula and only a few variables in the vv-list before this becomes
an overcrowded line, exceeding the line width and extending into and
perhaps beyond the margin. To split the display over two lines specify
a \verb`multline` or \verb`multline*` environment in the settings
option (if you don't want an equation number use the starred form):
\begin{verbatim}
    \solve[env=multline*,p=.]
      { ct-\sqrt{a^{2}+b^{2}-2ab\cos(\beta+\omega t)} 
      }[a=10,b=20,\beta=1,{t}=0][4]
\end{verbatim}
$\Longrightarrow$     \solve[env=multline*,p=.]
      { ct-\sqrt{a^{2}+b^{2}-2ab\cos(\beta+\omega t)} }
      [a=10,b=20,\beta=1,{t}=0][4]

\noindent For a formula with a significantly longer vv-list you could
introduce a third line to display the arrow and result on by entering
a specification like \verb`vv={,}\\(vv)\\` in the settings option
\emph{after} the \verb`env` setting.

In the example the function evaluates to $-0.0007$. Is this a zero
or an extremum? To find out, the calculation needs to be carried out
to a higher rounding value \textendash{} which is the reason why \verb`\nmcSolve`
has an extra rounding setting; see �\ref{subsec:solveExtraRounding}
below.

\subsection{\texttt{\textbackslash nmcSolve}-specific settings}

\begin{table}[b]
\centering{}\caption{\protect\label{tab:solveSettings}Settings for \texttt{\textbackslash nmcSolve}}
\begin{center}
\begin{tabular}{llll}
\toprule 
{\small key} & {\small type} & {\small meaning} & {\small initial}\tabularnewline
\midrule
{\small\texttt{var}} & {\small token(s)} & {\small equation variable} & \tabularnewline
{\small\texttt{dvar}} & {\small real $\ne0$} & {\small initial step size} & {\small\texttt{1}}\tabularnewline
{\small\texttt{+}} & {\small int} & {\small extra rounding} & {\small\texttt{0}}\tabularnewline
{\small\texttt{max}} & {\small int > 0} & {\small max. number of steps before cut off} & {\small\texttt{100}}\tabularnewline
\bottomrule
\end{tabular}
\par\end{center}
\end{table}
In addition to inherited settings there are some settings specific
to \verb`\nmcSolve`. These are listed in Table~\ref{tab:solveSettings}. 

\subsubsection{Equation variable}

By default the equation variable is the \emph{rightmost} variable
in the vv-list. This may not always be convenient. A different equation
variable can be specified by entering
\begin{verbatim}
    var = <variable name>
\end{verbatim}
in the vv-list. \verb`<variable name>` will generally be a single
character or token but multi-token names are perfectly acceptable
(with \texttt{numerica}'s default multi-token setting; see the associated
document \texttt{numerica.pdf} about this).

\subsubsection{Initial step size}

The vv-list contains the equation variable set to the initial trial
value. But \verb`\solve` needs \emph{two} initial values to begin
its search for a zero or extremum; see �\ref{subsec:solveSearch-strategy}.
Ideally, these values will straddle a zero of the function being investigated.
`Out of the box', the second trial value is $1$ more than the first:
if the equation variable is set to trial value $a$ then the second
value defaults to $a+1$. The `$+1$' here can be changed by entering
in the settings option
\begin{verbatim}
    dvar = <non-zero real number>
\end{verbatim}
For instance, \texttt{dvar=-1}, or \texttt{dvar=\textbackslash pi}
are two valid specifications of initial step size. The notation is
prompted by the use of expressions like $x$ and $x+dx$ for two nearby
points in calculus.

The initial step value may be too big or too small. An example where
the default step value is too big and a smaller one needs to be specified
is provided by Planck's radiation function (\emph{HMF }Table 27.2),
\[
f(x)=\frac{1}{x^{5}(e^{1/x}-1)}.
\]
From the (somewhat coarse-grained) table in \emph{HMF }it is clear
that there is a maximum of approximately 21.2 when $x$ is a little
more than $0.2$. This is a maximum above the $x$-axis and hence
`elusive' in the sense of �\ref{subsec:solveElusive-extrema}. To
find it, substract $100$ (say) from the formula and again use the
ability to nest commands to display the result. In the example, I
find in the vv-list of the \verb`\eval` command the value of $x$
which maximizes the Planck radiation function, then calculate the
maximum in the main argument of the \verb`\eval` command. Note the
\verb`dvar=0.1` in the settings option of the \verb`\solve` command: 
\begin{verbatim}
    \eval[p=.]{\[ \frac1{x^5(e^{1/x}-1)} \]}
      [ x={ \solve[dvar=0.1]
        { \frac1{y^5(e^{1/y}-1)}-100 }[y=0.1] } ]
\end{verbatim}
$\Longrightarrow$     \eval[p=.]{\[ \frac1{x^5(e^{1/x}-1)} \]}
      [ x={ \solve[dvar=0.1]
         { \frac1{y^5(e^{1/y}-1)}-100 }[y=0.1] } ]          

\noindent The maximum is indeed a little over $21.2$ and the $x$
value a little more than $0.2$. 

The default \verb`dvar=1` is too big for this problem. From the table
in \emph{HMF},\emph{ }$f(0.1)=4.540$ and $f(1.1)=0.419$. Thus for
$f(x)-100$ the `towards the $x$-axis' search strategy would lead
to negative values of $x$ with the default \verb`dvar` setting. 

\subsubsection{Extra rounding}

\label{subsec:solveExtraRounding}\verb`\solve` determines that a
zero or an extremum has been reached when the difference between two
successive bisection values vanishes at the specified rounding value
(the value in the final trailing optional argument of the \verb`\solve`
command; $6$ by default). If our function is $f(x)$ then $\abs{x_{n+1}-x_{n}}=0$
to the specified rounding value and $f(x_{n})$, $f(x_{n+1})$ have
opposite signs or at least one vanishes. Then (assuming $x_{n+1}>x_{n}$
and continuity) there must be a critical value $x_{c}\in[x_{n},x_{n+1}]$
such that $f(x_{c})=0$ exactly. But in general the critical value
$x_{c}$ will not coincide with $x_{n}$ or $x_{n+1}$. If $f(x)$
crosses the $x$-axis at a steep angle it may well be that although
$f(x_{c})$ vanishes to all $16$ figures, $f(x_{n})$ and $f(x_{n+1})$
do not, not even at the (generally smaller) specified rounding value.
For instance, suppose $f(x)=1000x-3000$ and that our trial value
is $x=e$:
\begin{centred}
\verb`\solve[vv=]{$ 1000x-3000 $}[x=e][4*]` $\Longrightarrow$ \solve[vv=]{$ 1000x-3000 $}[x=e][4*].
\end{centred}
Although the difference between successive $x$ values vanishes to
$4$ places of decimals, $f(x)$ does not, not even to $2$ places.
If we want the function to vanish at the specified rounding value
\textendash{} $4$ in the example \textendash{} then we will need
to locate the zero more precisely than that.

This is the purpose of the extra rounding key in the settings option.
Enter
\begin{verbatim}
    + = <integer>
\end{verbatim}
in the settings option of the \verb`\solve` command to add \verb`<integer>`
to the general rounding value. By default, \verb`+=0`. 

With this option available it is easy to check that \verb`+=3` suffices
in the example to ensure that both $x$ and $f(x)$ vanish to $4$
places of decimals,
\begin{centred}
\verb`\solve[+=3]{$ 1000x-3000 $}[x=e][4*]` $\Longrightarrow$ \solve[+=3]{$ 1000x-3000 $}[x=e][4*],
\end{centred}
and that \verb`+=2` does not, i.e., we need to locate the zero to
$4+3=7$ figures to ensure the function vanishes to $4$ figures. 

There is no need for the \verb`<integer>` to be positive. In fact
negative values can illuminate what is going on. In the first of the
following, the display is to $10$ places but (\verb`+=-4`) the calculation
is only to $10-4=6$ places. In the second, the display is again to
$10$ places, but (\verb`+=-3`) the calculation is to $10-3=7$ places.
\begin{centred}
\verb`\solve[+=-4]{$ 1000x-3000 $}[x=e][10*]` $\Longrightarrow$
\solve[+=-4]{$ 1000x-3000 $}[x=e][10*],

\verb`\solve[+=-3]{$ 1000x-3000 $}[x=e][10*]` $\Longrightarrow$
\solve[+=-3]{$ 1000x-3000 $}[x=e][10*].
\end{centred}
Only in the second does $f(x)=1000x-3000$ vanish when rounded to
$4$ figures. 

Returning to an earlier example (�\ref{subsec:solveMulti-line-display})
in which it was not entirely clear whether a zero or an extremum had
been found, we can now resolve the confusion. Use the extra rounding
setting (and pad with zeros to emphasize the $4$-figure display by
adding an asterisk in the trailing optional argument):
\begin{verbatim}
    \solve[env=multline*,p=.,+=2]
      { ct-\sqrt{a^{2}+b^{2}-2ab\cos(\beta+\omega t)} }
      [a=10,b=20,\beta=1,{t}=0][4*]
\end{verbatim}
$\Longrightarrow$     \solve[env=multline*,p=.,+=2]
      { ct-\sqrt{a^{2}+b^{2}-2ab\cos(\beta+\omega t)} }
      [a=10,b=20,\beta=1,{t}=0][4*]

\subsubsection{Maximum number of steps before cut-off}

Once two function values have been found of different sign, bisection
is guaranteed to arrive at a result. The problem is the \emph{search}
for two such values. This may not terminate \textendash{} think of
a function like $e^{x}$ which lacks both zeros and extrema. To prevent
an infinite loop, \verb`\solve` cuts off the search after $100$
steps. This cut-off value can be changed for a calculation by entering 
\begin{verbatim}
    max = <positive integer>
\end{verbatim}
in the settings option.

To illustrate, we know that $1/x$ has neither zero nor extremum,
but we do not get an infinite loop \textendash{} we get an error message
if we attempt to `solve' $1/x$:
\begin{centred}
\verb`\solve{ 1/x }[x=1]` $\Longrightarrow$ \solve{ 1/x }[x=1]
\end{centred}

\subsubsection{Form of result saved by \texttt{\textbackslash nmcReuse}}

\verb`\nmcReuse` saves (only) the numerical result. Version 2 offered
a setting \verb`reuse` providing a choice of what was saved. That
has been removed in version 3.

\chapter{Recurrence relations: \texttt{\textbackslash nmcRecur}}

One of the simplest recurrence relations is that determining the Fibonacci
numbers, $f_{n+2}=f_{n+1}+f_{n}$, with initial values $f_{0}=f_{1}=1$.
The command \verb`\nmcRecur`, short-name form \verb`\recur`, allows
calculation of the terms of this sequence:
\begin{verbatim}
    $ \nmcRecur[do=8,see1=8,...]
      { f_{n+2}=f_{n+1}+f_{n} }[f_{1}=1,f_{0}=1] $
\end{verbatim}
$\Longrightarrow$     $\nmcRecur[do=8,see1=8,...]
      { f_{n+2}=f_{n+1}+f_{n} }
        [f_{1}=1,f_{0}=1]$

The recurrence relation is entered in the main argument (between braces),
the initial values in the vv-list trailing the main argument, and
the display specification is placed in the settings option: \texttt{do=8}
terms to be calculated, all $8$ to be viewed (\texttt{see1=8}), and
the display to be concluded by an ellipsis to indicate that the sequence
continues (those are three dots/periods/full stops in the settings
option, not an ellipsis glyph).

A more complicated recurrence relation determines the Legendre polynomials:
\[
(n+2)P_{n+2}(x)-(2n+3)xP_{n+1}(x)+(n+1)P_{n}(x)=0.
\]
For the purposes of \verb`\recur` we need $P_{n+2}$ expressed in
terms of the lower order terms:
\[
P_{n+2}(x)=\frac{1}{n+2}\left((2n+3)xP_{n+1}(x)-(n+1)P_{n}(x)\right).
\]
 It is this standard form \textendash{} the term to be calculated
on the left, equated to an expression involving a fixed number of
lower-order terms on the right \textendash{} that \verb`numerica-plus`
works with. For $P_{0}(x)=1,~P_{1}(x)=x$ and $x=0.5$, the terms
are calculated like this:
\begin{verbatim}
    \recur[env=multline*,do=11,see1=4,see2=2,
      vv={,}\\(vv)\\] { P_{n+2}(x)=\frac{1}{n+2}
        \Bigl((2n+3)xP_{n+1}(x)-(n+1)P_{n}(x)\Bigr) }
      [P_{1}(x)=x,P_{0}(x)=1,x=0.5][7]
\end{verbatim}
$\Longrightarrow$     \recur[env=multline*,do=11,see1=4,see2=2,
      vv={,}\\(vv)\\,*]{ P_{n+2}(x)=\frac{1}{n+2}
        \Bigl((2n+3)xP_{n+1}(x)-(n+1)P_{n}(x)\Bigr) }
      [P_{1}(x)=x,P_{0}(x)=1,x=0.5][7] 

\noindent where $P_{9}(0.5)=-0.2678986$ and $P_{10}(0.5)=-0.1882286$
are the last two displayed values (and to seven figures are the values
listed in \emph{HMF }Table 8.1).

The examples also illustrate a consistent response with other commands
of the \texttt{numerica} suite to math environments. For the Fibonacci
sequence the \verb`\recur` command lay between \verb`$` delimiters
and display of the formula was suppressed. For the Legendre polynomials,
\verb`multline*` was specified in the settings option and the formula
was included in the display. (`Formula show', `formula hide' can
also be set with the \verb`f=1`, \verb`f=0` settings.) Note also
the specification of the vv-list in the second example, spreading
the display over three lines, and the \verb`see2` setting, specifying
how many terms to display. The first example displays a concluding
ellipsis; the second example doesn't. That is the result of the presence
of the \verb`...` (three periods) setting in the first example.

\section{Notational niceties}

More than any of the other commands in {\ttfamily\verb`numerica`}
and associated packages, \verb`\nmcRecur` depends on getting the
notation into a standard form. At least at this stage (version 3.0.0)
of \texttt{numerica-plus}
\begin{itemize}
\item the terms of the recurrence must be \emph{subscripted}: $f_{n}$,
$P_{n}(x)$ are examples;
\item the recurrence relation is placed in the main (mandatory) argument
of \verb`\nmcRecur` in the form: \emph{high-order term=function of
consecutive lower-order terms};
\item the initial-value terms in the vv-list must occur left-to-right in
the order \emph{high }to \emph{low} order;
\item the recurrence variable changes by $1$ between successive terms.
\end{itemize}
The example for Legendre polynomials in particular shows what is required.
The Fibonacci example is simpler, since the recurrence variable does
not occur independently in the recurrence relation as it does with
the Legendre polynomials, although in both cases the recurrence variable
is absent from the vv-list.

\subsection{Recurrence variable in the vv-list}

The recurrence variable is required in the vv-list only when an implicit
mode calculation is undertaken. To that end write $A$ and $B$ for
the coefficients $2n+3$ and $n+1$ respectively in the Legendre recurrence.
$A$ and $B$ will now need entries in the vv-list which means the
recurrence variable will need a value assigned to it there too, and
we will need to add \verb`vv@=1` (or \verb`vvmode=1`) to the settings
option.
\begin{verbatim}
    \recur[env=multline*,vvmode=1,do=11,see1=4,see2=2,
             ...,vv={,}\\(vv)\\]
      { P_{n+2}(x)=\frac{1}{n+2}
        \Bigl(AxP_{n+1}(x)-BP_{n}(x)\Bigr) }
      [P_{1}(x)=x,P_{0}(x)=1,x=0.5,A=2n+3,B=n+1,n=0]
\end{verbatim}
$\Longrightarrow$     \recur[env=multline*,vvmode=1,do=11,see1=4,see2=2,
             ...,vv={,}\\(vv)\\]
      { P_{n+2}(x)=\frac{1}{n+2}
        \Bigl(AxP_{n+1}(x)-BP_{n}(x)\Bigr) }
      [P_{1}(x)=x,P_{0}(x)=1,x=0.5,A=2n+3,B=n+1,n=0] 

Since the vv-list is evaluated from the right, the left-to-right high-to-low
ordering of the initial-value terms means the value of the lowest
order term is read first. Although \verb`numerica-plus` depends on
this order of occurrence of the terms, they do not need to be \emph{consecutive}
as in the examples so far (although it is natural to enter them in
this way). \texttt{numerica-plus} reads the value of the subscript
of only the right-most term (the lowest order term), increments it
by $1$ when reading the next recurrence term to the left, and so
on. The reading of the subscript of the lowest order term in the vv-list
provides the initial value of the recurrence variable.

In the following example I have placed other items between $P_{1}(x)$
and $P_{0}(x)$ in the vv-list (but maintained their left-to-right
order) and given the recurrence variable $n$ a ridiculous initial
value $\pi^{2}/12$. (Because of the order in which things get done
`behind the scenes', \emph{some} value is necessary so that the
$n$ in `$B=n+1$' does not generate an `unknown token' message.)
The result is unchanged.
\begin{verbatim}
    \recur[env=multline*,vvmode=1,do=11,see1=4,see2=2,
             ...,vv={,}\\(vv)\\]
      { P_{n+2}(x)=\frac{1}{n+2}
        \Bigl(AxP_{n+1}(x)-BP_{n}(x)\Bigr) }
      [A=2n+3,P_{1}(x)=x,B=n+1,n=\pi^2/12,P_{0}(x)=1,x=0.5]
\end{verbatim}
$\Longrightarrow$     \recur[env=multline*,vvmode=1,do=11,see1=4,see2=2,
             ...,vv={,}\\(vv)\\]
      { P_{n+2}(x)=\frac{1}{n+2}
        \Bigl(AxP_{n+1}(x)-BP_{n}(x)\Bigr) }
      [A=2n+3,P_{1}(x)=x,B=n+1,n=\pi^2/12,P_{0}(x)=1,x=0.5]

\subsection{Form of the recurrence relation}

As noted earler, the form of the recurrence must be entered in the
main argument in the form: \emph{highest order term = function of
consecutive lower order terms}. The number of lower\emph{ }order terms
is the order of the recurrence. The Fibonacci and Legendre polynomial
recurrences are both second order and presented in the form: \emph{term
$n+2$ = function of term $n+1$ and term $n$}. We could equally
have done
\begin{verbatim}
    \nmcRecur[p,do=8,see1=8,...]
      {$ f_{n}=f_{n-1}+f_{n-2} $}
        [f_{1}=1,f_{0}=1]
\end{verbatim}
$\Longrightarrow$     \nmcRecur[p,do=8,see1=8,...]
      {$ f_{n}=f_{n-1}+f_{n-2} $}
        [f_{1}=1,f_{0}=1] where now the recurrence is of the form \emph{term }$n$\emph{ =
function of term $n-1$ and term $n-2$}, or (adjusting the coefficients
as well as the recurrence terms),
\begin{verbatim}
    \recur[env=multline*,p=.,do=10,see1=4,see2=2,
             vv={,}\\(vv)\\]
      { P_{n+1}(x)=\frac{1}{n+1}
        \Bigl((2n+1)xP_{n}(x)-nP_{n-1}(x)\Bigr) }
          [P_{2}(x)=-0.125,P_{1}(x)=x,x=0.5]
\end{verbatim}
$\Longrightarrow$     \recur[env=multline*,p=.,do=10,see1=4,see2=2,
             vv={,}\\(vv)\\]
      { P_{n+1}(x)=\frac{1}{n+1}
        \Bigl((2n+1)xP_{n}(x)-nP_{n-1}(x)\Bigr) }
          [P_{2}(x)=-0.125,P_{1}(x)=x,x=0.5]

\noindent The recurrence here is of the form \emph{term }$n+1$\emph{
= function of term $n$ and term $n-1$}. This last example has one
further `wrinkle'. I've made $P_{1}(x)$ the lowest order term and
decreased the number of terms to calculate by $1$ accordingly. 

\subsection{First order recurrences (iteration)}

The recurrence relations for both the Fibonacci sequence and Legendre
polynomials are second order. There is no reason why the recurrence
should not be of third or higher order or, indeed, lower. A first
order recurrence provides an alternative means of iterating functions.
\verb`\recur` therefore provides a means to display the results of
an iteration in a different form from \verb`\iter`. 

Iterating $1+a/x$ in this way, $16$ terms gives the sequence
\begin{verbatim}
    \recur[do=16,see1=0,see2=3,...]{$
      x_{n+1}=1+a/x_{n}
        $}[x_{0}=1,a=1]
\end{verbatim}
$\Longrightarrow$     \recur[do=16,see1=0,see2=3,...]{$
      x_{n+1}=1+a/x_{n} 
        $}[x_{0}=1,a=1]

\noindent to be compared with the example near the start of Chapter~\ref{chap:iterIterating-functions}.
(\emph{That} effected $15$ iterations; \emph{this} uses $16$ terms
because of the extra $x_{0}=1$ term.) 

\section{Star (\texttt{{*}}) option}

When the star option is used with the \verb`\nmcRecur` command, only
a single term, the \emph{last}, is presented as the result. Repeating
the penultimate calculation, but with the star option produces
\begin{verbatim}
    \recur*[do=10]{ P_{n+1}(x)=\frac{1}{n+1}
        \Bigl((2n+1)xP_{n}(x)-nP_{n-1}(x)\Bigr) }
          [P_{2}(x)=-0.125,P_{1}(x)=x,x=0.5]
\end{verbatim}
$\Longrightarrow$     \recur*[do=10]{\[
      P_{n+1}(x)=\frac{1}{n+1}
        \Bigl((2n+1)xP_{n}(x)-nP_{n-1}(x)\Bigr)
          \]}[P_{2}(x)=-0.125,P_{1}(x)=x,x=0.5]

\noindent With the exception of \verb`do`, any settings are ignored
in the display of the result. The star option produces a purely numerical
answer without any trimmings.

\noindent{}%
\noindent\begin{minipage}[t]{1\columnwidth}%
\begin{shaded}%
This seems something of a waste of the star option since it gives
much the same result as choosing \texttt{do=10,see1=0,see2=1} \textendash{}
not \emph{exactly} the same, since math delimiters are involved here,
but sufficiently similar to make me wonder if I should change the
starred form to apply only to those recurrences which approach a limit.
The starred form would then produce the limiting value as its result
(like \verb`\iter*`). This is a possible change for future versions
of \texttt{numerica-plus} and should be borne in mind if using \verb`\recur*`.\end{shaded}%
\end{minipage}

\section{Settings}

The settings option is a comma-separated list of items of the form
\emph{key~=~value}. Because recurrence terms are necessarily multi-token,
the multi-token key is hard-coded in \verb`\recur` to \texttt{xx=1}. 

\subsubsection{Multi-line formatting of result}

When the \verb`\recur` command wraps around math delimiters, the
\texttt{vv} setting is available to split display of the result over
two or more lines. With the enhanced treatment of environments in
version 3 of \texttt{numerica}, some of which carries over to \texttt{numerica-plus},
it is possible to spread a display over two lines simply by specifying
\verb`env=multline*` (or \verb`env=multline` if you want equation
numbers). The default vv-list specification in this case is \verb`vv={,}\\(vv)`
which pushes the vv-list and sequence of calculated values to the
second line. If two lines still give a crowded result, \verb`vv={,}\\(vv)\\`
pushes the vv-list, centred, to a second line and the sequence of
values, right aligned, to a third. If you wanted only the sequence
of calculated values on a second line with the vv-list on the same
line as the formula then (for insance) \verb`vv={,}\qquad(vv)\\`
achieves this: 
\begin{verbatim}
    \nmcRecur[do=8,see1=8,...,vv={,}\qquad(vv)\\,*]
      {$ f_{n+2}=f_{n+1}+f_{n} $}
        [f_{1}=1,f_{0}=1]
\end{verbatim}
$\Longrightarrow$     \nmcRecur[do=8,see1=8,...,vv={,}\qquad(vv)\\,*]
      {$ f_{n+2}=f_{n+1}+f_{n} $}
        [f_{1}=1,f_{0}=1]

\subsection{\texttt{\textbackslash nmcRecur}-specific settings}

\label{subsec:recurSpecific-settings}
\begin{table}[b]
\centering{}\caption{\protect\label{tab:recurSettings}Settings for \texttt{\textbackslash nmcRecur}}
\begin{center}
\begin{tabular}{llll}
\toprule 
{\small key} & {\small type} & {\small meaning} & {\small default}\tabularnewline
\midrule
{\small\texttt{do}} & {\small int$\ge0$} & {\small number of terms to calculate} & {\small\texttt{7}}\tabularnewline
{\small\texttt{see1}} & {\small int$\ge0$} & {\small number of initial terms to display} & {\small\texttt{3}}\tabularnewline
{\small\texttt{see2}} & {\small int$\ge0$} & {\small number of final terms to display} & {\small\texttt{2}}\tabularnewline
{\small\texttt{...}} & {\small chars} & {\small follow display of values with an ellipsis} & \tabularnewline
\bottomrule
\end{tabular}
\par\end{center}
\end{table}
In addition to the inherited settings there are some specific to \verb`\nmcRecur`.
These are listed in Table~\ref{tab:recurSettings}.

\subsubsection{Number of terms to calculate}

By entering
\begin{verbatim}
    do = <integer>
\end{verbatim}
in the settings option you can specify how many terms of a recurrence
to calculate. The default is set to $7$ (largely to show a sufficient
number of terms of the Fibonacci series to begin to be interesting).
Note that \verb`<integer>` will generally \emph{not} correspond to
the subscript on the last term calculated since that also depends
on the value of the subscript of the lowest order term in the vv-list.

\subsubsection{Number of terms to display}

By entering
\begin{verbatim}
    see1 = <integer1>, see2=<integer2>
\end{verbatim}
in the settings option, you can specify how many initial terms of
the recurrence and how many of the final terms calculated you want
to view. If the sum of these settings is less than the \verb`do`
setting, then the terms are displayed with an intervening ellipsis.
If the sum is greater than the \verb`do` setting, then the values
are adjusted so that their sum equals the \verb`do` setting and all
terms are displayed. 

The adjustment is preferentially to \verb`see1`. Suppose \verb`do=7`,
\verb`see1=5`, \verb`see2=4`. Then \verb`see2` is left unchanged
but \verb`see1` is reduced to \verb`7-4=3`. If, say, \verb`do=7`,
\verb`see1=5`, \verb`see2=8`, then \verb`see2` is reduced to $7$
and \verb`see1` to \texttt{-1} (rather than zero, for technical reasons). 

The default value for \verb`see1` is $3$; the default value for
\verb`see2` is $2$.

\subsubsection{Ellipsis}

Including three dots (periods, fullstops) in the settings option 
\begin{verbatim}
    ...
\end{verbatim}
ensures that a (proper) ellipsis is inserted after the final term
is displayed. An example is provided by the display of the Fibonacci
sequence at the start of this chapter. By default this option is turned
off.

\subsection{Form of result saved by \texttt{\textbackslash nmcReuse}}

In previous versions of \texttt{numerica-plus} it was possible to
choose with the \verb`reuse` setting what was saved by the next \verb`\nmcReuse`
command. The \verb`reuse` setting has been discontinued in version
3 of \texttt{numeric-plus}. Now it is the last \verb`see2` values
of the recurrence sequence which are saved as a comma list in which
each saved value is wrapped in braces in case the decimal comma is
being used. (This setting has no effect when the star option is used
with \verb`\nmcRecur`. In that case only the numerical result of
the final term calculated is saved.) 

As an example, 
\begin{verbatim}
    \recur[env=multline*,p=.,vv@=1,do=11,see1=4,see2=2,
             vv={,}\\(vv)\\]
      { P_{n+2}(x)=\frac{1}{n+2} 
            \Bigl(kxP_{n+1}(x)-(n+1)P_{n}(x)\Bigr) }
      [k=2n+3,n=1,P_{1}(x)=x,P_{0}(x)=1,x=0.5]
    \reuse{legendre}
\end{verbatim}
$\Longrightarrow$     \recur[env=multline*,p=.,vv@=1,do=11,see1=4,see2=2,
             vv={,}\\(vv)\\ ]
      { P_{n+2}(x)=\frac{1}{n+2} 
            \Bigl(kxP_{n+1}(x)-(n+1)P_{n}(x)\Bigr) }
      [k=2n+3,n=1,P_{1}(x)=x,P_{0}(x)=1,x=0.5] \par
    \reuse[renew]{legendre}

\noindent Now check to see what has been saved: 
\begin{centred}
\verb`$\legendre$` $\Longrightarrow$  $ \legendre $.
\end{centred}
As you can see, the final two (because of \verb`see2=2`) of the $11$
Legendre polynomials calculated ($P_{0}(x)$ is the first) have been
saved. 

\subsubsection{Accessing individual terms}

\label{subsec:recurAccessing-individual-terms}You may wish to gain
access to individual members of the sequence of saved values. \texttt{numerica-plus}
provides the utility function \verb`\clitem` for this purpose.\footnote{In fact a wrapper around the \texttt{expl3} function \texttt{\textbackslash clist\_item:Nn}.}
It acts on the macro containing the comma list and a number specifying
which member of the sequence is to be recovered: \verb`1` for the
first (leftmost) item, \verb`2` for the second item, and so on. 
\begin{centred}
\verb`$\clitem\legendre 2$` $\Longrightarrow$ $\clitem\legendre 2$
\end{centred}
The \verb`$` delimiters ensure the minus sign displays correctly.
Multi-digit numbers need to be enclosed in braces. Negative integers
give access to the end of the sequence, \verb`-1` for the last (rightmost)
item, \verb`-2` for the second last and so on. In the present case,
\begin{centred}
\verb`$\clitem\legendre{-1}=\clitem\legendre 2$` $\Longrightarrow$
$\clitem\legendre{-1}=\clitem\legendre 2$.
\end{centred}

\subsection{Orthogonal polynomials}

I've used Legendre polynomials in examples above, but orthogonal polynomials
generally lend themselves to the \verb`\recur` treatment. Quoting
from \emph{HMF} 22.7, orthogonal polynomials $f_{n}$ satisfy recurrence
relations of the form
\[
a_{1n}f_{n+1}(x)=(a_{2n}+a_{3n}x)f_{n}(x)-a_{4n}f_{n-1}(x),
\]
or in the standard form required by \verb`\recur`,
\[
f_{n+1}(x)=\frac{a_{2n}+a_{3n}x}{a_{1n}}f_{n}(x)-\frac{a_{4n}}{a_{1n}}f_{n-1}(x).
\]
\emph{HMF} 22.7 provides a listing of the coefficients $a_{in}$ for
the polynomials of Jacobi, Chebyshev, Legendre, Laguerre, Hermite
and others, and tables for these polynomials.

For example, Laguerre polynomials satisfy the recurrence 
\[
L_{n+1}(x)=\frac{2n+1-x}{n+1}L_{n}(x)-\frac{n}{n+1}L_{n-1}(x).
\]
with initial values $L_{0}(x)=1$ and $L_{1}(x)=1-x$. So let's calculate
the first $13$ Laguerre polynomials for, say, $x=0.5$: 
\begin{verbatim}
    \recur[env=multline*,do=13,see1=4,see2=2,
             vv={,}\\(vv)\\]
      { L_{n+1}(x)=\frac{2n+1-x}{n+1}L_{n}(x)-
        \frac{n}{n+1}L_{n-1}(x) }
           [L_{1}(x)=1-x,L_{0}(x)=1,x=0.5]
\end{verbatim}
$\Longrightarrow$     \recur[env=multline*,do=13,see1=4,see2=2,
             vv={,}\\(vv)\\]
      { L_{n+1}(x)=\frac{2n+1-x}{n+1}L_{n}(x)-
        \frac{n}{n+1}L_{n-1}(x) }
          [L_{1}(x)=1-x,L_{0}(x)=1,x=0.5]

\noindent and for $x=5$: 
\begin{verbatim}
    \recur[env=multline*,p=.,do=13,see1=4,see2=2,
             vv={,}\\(vv)\\]
      { L_{n+1}(x)=\frac{2n+1-x}{n+1}L_{n}(x)-
        \frac{n}{n+1}L_{n-1}(x) }
          [L_{1}(x)=1-x,L_{0}(x)=1,x=5]
\end{verbatim}
$\Longrightarrow$     \recur[env=multline*,p=.,do=13,see1=4,see2=2,
             vv={,}\\(vv)\\]
      { L_{n+1}(x)=\frac{2n+1-x}{n+1}L_{n}(x)-
        \frac{n}{n+1}L_{n-1}(x) }
          [L_{1}(x)=1-x,L_{0}(x)=1,x=5]

\noindent The results (reassuringly) coincide with those provided
in \emph{HMF }Table 22.11.

\subsection{Nesting}

It is possible to use the \verb`\recur` command within an \verb`\eval`,
\verb`\iter`, or \verb`\solve` command, and indeed in \verb`\recur`
itself, but with this caveat: if \verb`\recur` is nested within another
command, the initial terms of the recurrence \textendash{} e.g., $f_{1}=1,f_{0}=1$,
for the Fibonacci series, or $L_{1}(x)=1-x,L_{0}(x)=1$ for the Laguerre
polynomials \textendash{} must be located in the vv-list of that\emph{
inner }\verb`\recur`\emph{ }command. Other shared variables can often
be shifted to the vv-list of the outer command, but not these initial
terms.

In the following example I multiply together (rather futilely) the
third and fourth members of the sequence of Laguerre polynomials for
$x=5$ (the answer expected is \verb`$ \eval{3.5\times2.666667} $`
$\Longrightarrow$ $ \eval{3.5\times2.666667} $). Note that although
it is tempting to shift the shared vv-lists of the inner \verb`\recur*`
commands to the vv-list of the outer \verb`\eval` command, in fact
only the \verb`x=5` entry has been transferred:
\begin{verbatim}
    \eval[p=.]{$
      \recur[do=3]
        { L_{n+1}(x)=\frac{2n+1-x}{n+1}L_{n}(x)-
            \frac{n}{n+1}L_{n-1}(x)}
          [L_{1}(x)=1-x,L_{0}(x)=1]
      \times 
      \recur[do=4]
        { L_{n+1}(x)=\frac{2n+1-x}{n+1}L_{n}(x)-
            \frac{n}{n+1}L_{n-1}(x)}
          [L_{1}(x)=1-x,L_{0}(x)=1]
    $}[x=5]
\end{verbatim}
$\Longrightarrow$     \eval[p=.]{$
      \recur[do=3]
        { L_{n+1}(x)=\frac{2n+1-x}{n+1}L_{n}(x)-
            \frac{n}{n+1}L_{n-1}(x)}
              [L_{1}(x)=1-x,L_{0}(x)=1,x=5]
      \times 
      \recur[do=4]
        { L_{n+1}(x)=\frac{2n+1-x}{n+1}L_{n}(x)-
            \frac{n}{n+1}L_{n-1}(x)}
              [L_{1}(x)=1-x,L_{0}(x)=1,x=5]
    $}

\noindent{}%
\noindent\begin{minipage}[t]{1\columnwidth}%
\begin{shaded}%
The terms of a recurrence relation are multi-token variables but \texttt{numerica}
requires single tokens for its calculations. The problem for \verb`\recur`
is that the terms in the recurrence relation in the main (mandatory)
argument differ from the terms in the vv-list: for instance $f_{n}$
in the main argument, $f_{0}$ in the vv-list. If left like that,
when \texttt{numerica} does its conversion from multi-token to single
token variables, $f_{n}$ would not be found since it differs from
$f_{0}$. Hence a crucial first step for \verb`\recur` is to reconcile
the different forms, which it does by converting the forms in the
vv-list to the forms in the recurrence in the main argument. To be
available for this form change, they must reside in the \emph{inner}
vv-list. In the outer vv-list they would be inaccessible to the inner
command.

{*}{*}{*}

This suggests an alternative way of proceeding: write the inital values
of the recurrence terms in the \emph{same} form in which they occur
in the recurrence relation, together with an initial value for the
recurrence variable, e.g., $f_{n+1}=1,f_{n}=1,n=0$. This is not how
mathematicians write the initial values in recurrence relations, which
is why I did not pursue it, but it neatly sidesteps what is otherwise
an initial awkwardness. \end{shaded}%
\end{minipage}

\chapter{Reference summary}

\section{Commands defined in \texttt{numerica-plus}}
\begin{enumerate}
\item \texttt{\textbackslash nmcIterate, \textbackslash iter}
\item \texttt{\textbackslash nmcSolve, \textbackslash solve}
\item \textbackslash\texttt{nmcRecur, \textbackslash recur}
\item \texttt{\textbackslash clitem}
\end{enumerate}

\section{Settings for the three main commands}

\subsection{Settings for \texttt{\textbackslash nmcIterate}}

Settings keys for \verb`\nmcIterate`:

\begin{center}
\begin{tabular}{llll}
\toprule 
{\small key} & {\small type} & {\small meaning} & {\small initial}\tabularnewline
\midrule
{\small\texttt{var}} & {\small token(s)} & {\small iteration variable} & \tabularnewline
{\small\texttt{+}} & {\small int} & {\small fixed point extra rounding} & {\small\texttt{0}}\tabularnewline
{\small\texttt{max}} & {\small int > 0} & {\small max. iteration count (fixed points)} & {\small\texttt{100}}\tabularnewline
{\small\texttt{do}} & {\small int > 0} & {\small number of iterations to perform} & {\small\texttt{5}}\tabularnewline
{\small\texttt{see}} & {\small int > 0} & {\small number of final iterations to view} & {\small\texttt{4}}\tabularnewline
\bottomrule
\end{tabular}
\par\end{center}

\subsection{Settings for \texttt{\textbackslash nmcSolve}}

Settings keys for \verb`\nmcSolve`:

\begin{center}
\begin{tabular}{llll}
\toprule 
{\small key} & {\small type} & {\small meaning} & {\small initial}\tabularnewline
\midrule
{\small\texttt{var}} & {\small token(s)} & {\small equation variable} & \tabularnewline
{\small\texttt{dvar}} & {\small real $\ne0$} & {\small initial step size} & {\small\texttt{1}}\tabularnewline
{\small\texttt{+}} & {\small int} & {\small extra rounding} & {\small\texttt{0}}\tabularnewline
{\small\texttt{max}} & {\small int > 0} & {\small max. number of steps before cut off} & {\small\texttt{100}}\tabularnewline
\bottomrule
\end{tabular}
\par\end{center}

\subsection{Settings for \texttt{\textbackslash nmcRecur}}

Settings keys for \verb`\nmcRecur`:

\begin{center}
\begin{tabular}{llll}
\toprule 
{\small key} & {\small type} & {\small meaning} & {\small initial}\tabularnewline
\midrule
{\small\texttt{do}} & {\small int$\ge0$} & {\small number of terms to calculate} & {\small\texttt{7}}\tabularnewline
{\small\texttt{see1}} & {\small int$\ge0$} & {\small number of initial terms to display} & {\small\texttt{3}}\tabularnewline
{\small\texttt{see2}} & {\small int$\ge0$} & {\small number of final terms to display} & {\small\texttt{2}}\tabularnewline
{\small\texttt{...}} & {\small chars} & {\small follow display of values with an ellipsis} & \tabularnewline
\bottomrule
\end{tabular}
\par\end{center}
\end{document}
